\documentclass[11pt,a4paper]{article}
\input{../../_archive/tfpt-45/tfpt45-common}

% Override the running header inherited from tfpt45-common.tex
% to match this note's actual topic, not the parent paper series.
\fancyhf{}
\fancyhead[L]{\small\itshape\color{tfptgray} TFPT Lean Note: Carrier Rigidity Audit Harness}
\fancyhead[R]{\small\color{tfptgray}\thepage}
\renewcommand{\headrulewidth}{0.4pt}

\usepackage{listings}
\usepackage{xcolor}

\definecolor{leanbg}{RGB}{248,248,250}
\definecolor{leankeyword}{RGB}{35,60,120}
\definecolor{leancomment}{RGB}{90,90,95}
\definecolor{leanstring}{RGB}{40,110,60}

\lstdefinelanguage{lean4}{
  morekeywords={import,namespace,end,theorem,lemma,def,noncomputable,structure,
    where,by,have,show,rw,unfold,calc,obtain,refine,exact,fun,if,then,else,
    match,with,let,in,return,do,abbrev,instance,open,variable,inductive,class},
  sensitive=true,
  morecomment=[s][\color{leancomment}\itshape]{/-}{-/},
  morecomment=[l][\color{leancomment}\itshape]{--},
  morestring=[b]",
}

\lstset{
  language=lean4,
  basicstyle=\ttfamily\footnotesize,
  keywordstyle=\color{leankeyword}\bfseries,
  commentstyle=\color{leancomment}\itshape,
  stringstyle=\color{leanstring},
  backgroundcolor=\color{leanbg},
  frame=single,
  framerule=0.4pt,
  rulecolor=\color{tfptblue!40},
  breaklines=true,
  showstringspaces=false,
  columns=fullflexible,
  xleftmargin=4pt, xrightmargin=4pt,
  aboveskip=6pt, belowskip=6pt,
  literate={λ}{{$\lambda$}}1 {α}{{$\alpha$}}1 {β}{{$\beta$}}1
           {ε}{{$\varepsilon$}}1 {ℚ}{{$\mathbb Q$}}1 {ℤ}{{$\mathbb Z$}}1
           {ℕ}{{$\mathbb N$}}1 {₊}{$_{+}$}1 {₋}{$_{-}$}1 {₀}{$_0$}1
           {⊕}{{$\oplus$}}1 {⟨}{{$\langle$}}1 {⟩}{{$\rangle$}}1
           {→}{{$\rightarrow$}}1 {↦}{{$\mapsto$}}1 {∀}{{$\forall$}}1
           {∃}{{$\exists$}}1 {·}{{$\cdot$}}1 {≠}{{$\neq$}}1
           {≤}{{$\leq$}}1 {≥}{{$\geq$}}1 {∈}{{$\in$}}1 {𝒫}{{$\mathcal P$}}1
           {∣}{{$\,|\,$}}1 {∘}{{$\circ$}}1 {ₗ}{$_{\ell}$}1
           {⅟}{{$\,\widehat{\cdot}\,$}}1 {≃}{{$\simeq$}}1 {Λ}{{$\Lambda$}}1
}

\title{\color{tfptblue}\textbf{Formal Verification of TFPT Carrier Rigidity in Lean 4}\\[0.4em]
{\Large An Audited Dependency Graph from the Boundary Involution\\
to the Standard-Model \boldmath$3+2$\unboldmath\ Hypercharge Carrier}}
\author{Stefan Hamann \and Alessandro Rizzo}
\date{{\normalsize\color{tfptgray}\sffamily TFPT Experiment Note --- Lean 4 Formalisation --- v11 --- \today}}

\begin{document}
\maketitle

\begin{abstract}
This note documents a Lean~4 formalisation of the algebraic
carrier rigidity theorem of TFPT Paper~2 (\emph{Carrier Rigidity,
Hypercharge, and the Standard-Model Packet from Boundary
Polarization}), up to and including the determinant-normalised
$3+2$ hypercharge carrier
$Y = \operatorname{diag}(-\tfrac13,-\tfrac13,-\tfrac13,\tfrac12,\tfrac12)$.
The bundled top-level theorem consumes four explicit interfaces:
a carrier involution, Higgs and Yukawa rank witnesses, and a
primitive oriented determinant-preserving carrier condition.
Each interface is refinable upstream via a separately formalised
Lean theorem. In v10 each upstream interface is closed by a
further structural layer: the Calder\'on idempotent
$\pi^2 = \pi$ is itself produced from a
\texttt{BoundaryPolarization} ($H = H_+ \oplus H_-$); the Higgs
algebraic shadow is identified with the degree-1 piece of the
bivariate polynomial ring
$K[X_0, X_1]_1 = \texttt{MvPolynomial.homogeneousSubmodule (Fin 2) K 1}$
(the algebraic-geometry "Proj-side" presentation of
$H^0(\mathbb P^1, \mathcal O(1))$); and the Stage D Yukawa
trilinear form is constructed explicitly from any basis of a
rank-$3$ space, with injective and surjective contraction proved
by direct basis computation. The Yukawa Stage D characterisation
$\exists\,\omega\;(\textrm{CI}\,\wedge\,\textrm{CS})
  \iff \dim E = 3$ now holds under \texttt{Nontrivial}\,$E$. The unique integer
pair $(q_-,q_+)=(-2,3)$ follows from a general $(m,n)$
lattice-rigidity theorem applied to the structurally bundled
\texttt{PrimitiveOrientedDeterminantCarrier 3 2} input. Lean
proves the unique $3+2$ hypercharge carrier, trace zero, and the
carrier polynomial $6 Y^2 - Y - \mathbf 1 = \mathbf 0$. The
project builds against Mathlib~v4.29.1 with zero \texttt{sorry},
zero \texttt{admit}, no domain-specific axioms, no
\texttt{unsafe}/\texttt{opaque}/\texttt{partial} escape hatches,
and every load-bearing theorem depends only on the three standard
Lean axioms (\texttt{propext}, \texttt{Classical.choice},
\texttt{Quot.sound}). A seven-check \texttt{audit.sh} gate is
shipped; the headline-theorem checks combine Lean
\texttt{\#check} elaboration (\texttt{AuditCheck.lean}) with
strict \texttt{example : <exact type> := <theorem>} signature
locks (\texttt{AuditContract.lean}), so neither silent renames
nor silent weakening of theorem statements can mask regressions.
\end{abstract}

\begin{tfptfrontbox}
\textbf{Inputs from previous papers.} Paper~1 (boundary primitive
kernel) for the \emph{conceptual} derivation of the carrier
involution from the Calder\'on polarisation. The Lean project
consumes only algebraic data: an involution element $\varepsilon$ with
$\varepsilon^2=\mathbf 1$, plus Stage-A certificates for the Higgs and
Yukawa rank discharges.

\textbf{New contribution.} A reproducible Lean~4 project
\texttt{TfptCarrier} (Mathlib~v4.29.1) that elevates four
previously imported inputs to formal Lean theorems:

\begin{itemize}[leftmargin=1.5em, itemsep=0pt]
    \item boundary involution $\Rightarrow$ orthogonal idempotents;
    \item general integer rigidity $\Rightarrow$ unique primitive pair;
    \item structural trace identity $\Rightarrow$ $\tr Y=0$ from
          rank data alone;
    \item determinant character of the carrier torus
          $\Rightarrow$ Diophantine condition $m q_- + n q_+ = 0$
          as a corollary (universal-$\lambda$ form, with a single-
          point helper at $\lambda=2$).
\end{itemize}

Two further inputs ($\dim E_+=2$ and $\dim E_-=3$) are represented
as structured hypotheses, not bare numerical assumptions. They
mark the exact interfaces where the future geometric (Higgs) and
representation-theoretic (Yukawa) proofs must enter.

For the Yukawa case, this note now provides \emph{four}
interface layers, with the upstream stage carrying genuine
algebraic content:

\begin{itemize}[leftmargin=1.5em, itemsep=1pt]
\item \emph{Stage A} (\texttt{YukawaRankCertificate}) --- a
    bare linear equivalence $E_- \simeq K^3$. Used by the
    bundled \texttt{CarrierData.CarrierPremises} as a rank
    witness; refinable from any of the upstream stages.
\item \emph{Stage B} (\texttt{YukawaTopForm}) --- a linear
    equivalence $\Lambda^3 E_- \simeq K$ (an algebraic top
    form on $E_-$).
\item \emph{Stage C} (\texttt{PrimitiveIndecomposableYukawaCoupling})
    --- a linear equivalence $E_- \simeq (\Lambda^2 E_-)^*$
    plus non-triviality (the contraction map of a non-degenerate
    alternating trilinear coupling, taken as data).
\item \emph{Stage D} (\texttt{YukawaTrilinearForm}) --- introduced
    in v6, with operational predicates renamed in v7 --- an
    alternating trilinear form $\omega : \Lambda^3 E_- \to K$ as
    primary datum, together with two explicit \texttt{Prop}
    predicates \texttt{ContractionInjective} and
    \texttt{ContractionSurjective}. From these the contraction
    map is \emph{constructed} (via
    \texttt{alternatingMapLinearEquiv} and
    \texttt{AlternatingMap.curryLeft}) and proved to be a linear
    iso when both predicates hold and $E_-$ is non-trivial,
    producing Stage C as a \emph{theorem}, not as a definition.
\end{itemize}

Each upstream stage produces the next via a Lean bridge
(D $\Rightarrow$ C $\Rightarrow$ B $\Rightarrow$ A). Stage D
carries the genuine content of TFPT Paper 2's
"non-degenerate primitive alternating trilinear coupling": the
trilinear form is the primary datum, the contraction iso is a
derived fact, and the rank-3 conclusion follows from the
contraction iso via the elementary dimension chain
$\dim E_- = \dim (\Lambda^2 E_-)^* = \binom{\dim E_-}{2}$ on a
non-trivial finite-dimensional space.

A seven-check \texttt{audit.sh} CI gate is included.

\textbf{Not claimed here.} No analytic content of Paper~1, no full
geometric formalisation of $H^0(\mathbb P^1, \mathcal O(1))$, no
quiver-theoretic definition of \emph{primitive indecomposable
Yukawa type}, no SM hypercharge spectrum on derived
representations (duals, exterior powers, tensor products), no
$\mathbb Z_6$ quotient, no $\alpha$, no flavour, no cosmology.

\textbf{Falsification or audit surface.} The note fails if
(i)~\texttt{lake build} no longer succeeds against the pinned
Mathlib release; (ii)~the \texttt{audit.sh} script returns a
non-zero exit; (iii)~the dependency graph stated below drifts from
the wording of Paper~2~\S2.
\end{tfptfrontbox}

\begin{tfptclaimcontract}
\textbf{Claim.}
The top-level Lean theorem
\texttt{CarrierPremises.hypercharge\_carrier\_packet} takes
four explicit interfaces:

\begin{enumerate}[leftmargin=1.5em, itemsep=0pt, label=(\arabic*)]
\item a carrier involution element $\varepsilon$ with
      $\varepsilon^2 = \mathbf 1$
      (\texttt{InvolutionProjectors.CarrierInvolution}, optionally
      wrapped by a \texttt{CalderonCertificate});
\item a Higgs rank-2 certificate
      (\texttt{HiggsIndexCertificate}) discharging $\dim E_+ = 2$;
\item a Yukawa rank-3 witness
      (\texttt{YukawaRankCertificate}) discharging $\dim E_- = 3$,
      \emph{refinable} via \texttt{YukawaTopForm} (Stage B),
      \texttt{PrimitiveIndecomposableYukawaCoupling} (Stage C),
      and \texttt{YukawaTrilinearForm} (Stage D) by formal Lean
      bridges;
\item a primitive oriented determinant-preserving carrier
      condition
      $(3 q_- + 2 q_+ = 0,\ q_- < 0 < q_+,\ \gcd(|q_-|,q_+) = 1)$,
      packaged as
      \texttt{PrimitiveOrientedDeterminantCarrier 3 2}.
\end{enumerate}

\textbf{Conclusion.}
From these inputs Lean derives the $3+2$ hypercharge carrier
$Y = -\tfrac13 P_- + \tfrac12 P_+$, with $(q_-,q_+) = (-2, 3)$,
$\tr Y = 0$, and the carrier polynomial
$6 Y^2 - Y - \mathbf 1 = \mathbf 0$ in $\mathrm{End}_K M$.

\textbf{Three layers of the contract.}
The claim is structured into three explicit layers, separated to
prevent overclaiming and to make every open question visible:

\smallskip
\noindent
\emph{(a) Direct inputs of \texttt{CarrierPremises}.}\ The
bundled top-level theorem consumes the four interfaces above:
a carrier involution, a Higgs rank witness, a Yukawa rank
witness, and a primitive oriented determinant-preserving
carrier condition.

\smallskip
\noindent
\emph{(b) Lean-proved upstream refinements (closed to the
next algebraic layer).}\ Each of the four direct inputs is
produced by an explicit Lean-proved upstream interface:

\begin{itemize}[leftmargin=1.6em, itemsep=0pt]
  \item Carrier involution $\leftarrow$
    \texttt{CalderonProjector} $\leftarrow$
    \texttt{BoundaryPolarization} ($H = H_+ \oplus H_-$).
    The converse \texttt{BoundaryPolarization.ofCarrierInvolution}
    makes the chain a Lean-proved \emph{algebraic
    equivalence}: the carrier involution and the polarisation
    are two presentations of the same finite-dimensional
    carrier datum. This is not the analytic Paper-1 content;
    the analytic construction of the polarisation from
    boundary elliptic theory remains upstream.
  \item Higgs rank witness $\leftarrow$
    \texttt{HiggsTopForm} (Stage B) $\leftarrow$
    \texttt{HiggsSchemeCohomologyShadow}
    ($\simeq \texttt{MvPolynomial.homogeneousSubmodule (Fin 2) K 1}$).
  \item Yukawa rank witness $\leftarrow$
    \texttt{YukawaTopForm} (Stage B) $\leftarrow$
    \texttt{YukawaPrimitive} (Stage C) $\leftarrow$
    \texttt{YukawaTrilinearForm} (Stage D) $\leftrightarrow$
    \texttt{YukawaStageDExistence} (\emph{both directions
    formalised, under} \texttt{Nontrivial}).
  \item Primitive oriented determinant carrier $\leftarrow$
    \texttt{SeamWindingInterface} (algebraic packaging of
    seam-class data).
\end{itemize}

\smallskip
\noindent
\emph{(c) Open TFPT-origin theorems (not formalised here).}\ Three
genuinely upstream questions remain open:

\begin{enumerate}[leftmargin=1.6em, itemsep=0pt, label=(\arabic*)]
  \item \emph{Paper 1 analytic boundary theory} produces the
    polarisation $H = H_+ \oplus H_-$ (boundary elliptic
    theory, Sobolev spaces, Calder\'on operator).
  \item \emph{TFPT boundary kernel} canonically produces an
    $\omega : \Lambda^3 E_- \to K$ with
    \texttt{ContractionInjective} and
    \texttt{ContractionSurjective} (the \emph{canonical}
    question; the existence is now closed by the present
    note).
  \item \emph{Seam winding $[u_\Sigma] = 1$} produces a
    primitive oriented determinant-preserving carrier
    condition.
\end{enumerate}

These three are \emph{outside the scope} of the present Lean
note; each is named, typed, and provided with a Lean-side
target structure
(\texttt{BoundaryPolarization},
\texttt{BoundaryYukawaKernelInterface},
\texttt{SeamWindingInterface}).

\textbf{Not claimed.}
Beyond the three open TFPT-origin theorems, the
$\mathbb Z_6$ quotient $G_{\mathrm{phys}}$, derived SM
representations (duals, exterior powers, tensor products),
$\alpha$, flavour, and cosmology are also not formalised
in this note.

\textbf{First assumptions on $K$.}
The projector construction requires $2$ invertible in the
carrier algebra; for vector-space applications this is ensured
by $\mathrm{char}\,K \ne 2$. The rational normalisation
$Y = -P_-/3 + P_+/2$ further requires $6$ invertible; concretely
the bundled theorems use \texttt{CharZero K}.

\textbf{Proof status.}
\emph{Formally verified} in Lean~4 + Mathlib~v4.29.1; no
\texttt{sorry}; no \texttt{admit}; no
\texttt{unsafe}/\texttt{opaque}/\texttt{partial}; no
domain-specific axiom; \texttt{audit.sh} (seven checks
including the signature-lock \texttt{AuditContract.lean})
returns \texttt{AUDIT:\ PASS}. As of v10, the three upstream
interfaces (Calder\'on, Higgs, Yukawa Stage D backward) are
each closed \emph{to the next algebraic layer}. The remaining
gaps are TFPT-side or analytic, not Lean-side.

\textbf{Kill condition.}
Any of the seven audit checks fail; any headline theorem in
\texttt{AuditCheck.lean} is removed or fails to elaborate;
any \texttt{example} in \texttt{AuditContract.lean} fails to
typecheck (which detects silent weakening of theorem statements).
\end{tfptclaimcontract}

\tableofcontents

% ---------------------------------------------------------------- %
\section{What This Note Is About}
% ---------------------------------------------------------------- %

v1 formalised the \emph{algebraic motor} of TFPT Paper~2 ---
the carrier polynomial and the trace identity on a concrete
$5\times5$ diagonal model. v2 began making the upstream
dependencies explicit. v3 made the dependency graph correct.
v4 replaced the Yukawa rank witness with a structurally
algebraic Stage B (top form on $\Lambda^3 E_-$). v5 introduced
a Stage C contraction iso $E_- \simeq (\Lambda^2 E_-)^*$.

v6 completed the Yukawa refinement chain by introducing a
Stage D primary datum; v7 renamed its predicates to operational
names; v8 added a Calder\'on idempotent layer and a Higgs Stage B;
v9 consolidates the document into a single review-ready state
and adds an unconditional transparency theorem for Stage D.
The Stage D primary datum reads:

\[
\texttt{YukawaTrilinearForm}\,K\,E_-
\;\;:=\;\;
\bigl\{\,\omega : \Lambda^3 E_- \to_{\ell K} K\,\bigr\},
\qquad
\texttt{ContractionInjective}\,\omega,\;\;
\texttt{ContractionSurjective}\,\omega.
\]

The contraction map $C_\omega : E_- \to (\Lambda^2 E_-)^*$ is
now \emph{constructed} from $\omega$ via
\texttt{exteriorPower.alternatingMapLinearEquiv} composed with
\texttt{AlternatingMap.curryLeft}. The two predicates state
exactly that $C_\omega$ is injective (mathematical reading:
non-degeneracy) and surjective (mathematical reading: primitivity
/ no spectator). Together, on a non-trivial finite-dimensional
space, they yield a Stage C structure as a theorem; the rank-3
conclusion then follows mechanically:

\[
\boxed{\;
\omega\;\text{trilinear} \;+\; \texttt{CI} \;+\; \texttt{CS} \;+\; \texttt{Nontrivial}\,E_-
\;\;\xRightarrow{\text{contraction is iso}}\;\;
\text{Stage C}
\;\;\xRightarrow{\binom{n}{2} = n,\, n\ne 0}\;\;
\dim E_- = 3.
\;}
\]

\noindent
The non-triviality hypothesis is essential: for $E_- = 0$ both
predicates are vacuously satisfied (the contraction $0 \to 0$ is
trivially bijective) but $\dim E_- = 0$. The Lean theorem
\texttt{finrank\_E\_eq\_zero\_or\_three} (new in v9) records the
unconditional disjunction $\dim E_- = 0 \lor \dim E_- = 3$ for
transparency; the rank-3 slogan above is the case under
\texttt{Nontrivial}\,$E_-$.

The full Yukawa tower in Lean is now
\[
\text{Stage D}
\;\Rightarrow\; \text{Stage C}
\;\Rightarrow\; \text{Stage B}
\;\Rightarrow\; \text{Stage A}
\;\Rightarrow\; \dim E_- = 3,
\]
with every step a Lean theorem and Stages B and C carrying
characterisation theorems establishing equivalence to
$\dim E_- = 3$.

v6 made the Paper~1 input explicit. A new
\texttt{CalderonInterface.lean} introduces a
\texttt{CalderonCertificate} structure that names the
Calder\'on-induced carrier involution as a Lean datum and feeds
it into the \texttt{InvolutionProjectors} chain. v8 upgrades
this interface from documentary to a structural algebraic step:
a new \texttt{CalderonProjector.lean} module shows that an
idempotent element $\pi$ in a ring with $2$ invertible yields
$\varepsilon^2 = 1$ via the algebraic identity $\varepsilon := 2\pi - 1$,
$(2\pi - 1)^2 = 4\pi^2 - 4\pi + 1 = 1$. v10 adds the next
upstream layer: a new \texttt{BoundaryPolarization.lean} module
shows that any direct-sum decomposition $H = H_+ \oplus H_-$
of a $K$-module produces an idempotent endomorphism on $H$ (via
\texttt{Submodule.IsCompl.projection}) and feeds the Calder\'on
chain. The remaining Paper-1 gap is therefore precisely the
analytic content: producing the boundary spectral
decomposition itself.

v8 also introduces a Stage-B Higgs interface
(\texttt{HiggsTopForm.lean}) parallel to the Yukawa tower:
the rank-2 conclusion $\dim E_+ = 2$ is derivable from a
structural hypothesis identifying $E_+$ with the algebraic
shadow of $H^0(\mathbb P^1, \mathcal O(1))$, namely the space
of homogeneous degree-1 polynomials in two formal variables
$\alpha X + \beta Y$. v10 closes the bridge to actual algebraic
geometry: \texttt{HiggsSchemeCohomologyShadow.lean} provides a
\texttt{LinearEquiv} between this algebraic shadow and the
Mathlib-native degree-1 piece of the bivariate polynomial ring
$\texttt{MvPolynomial.homogeneousSubmodule (Fin 2) K 1}
= K[X_0, X_1]_1$, which by the standard $\operatorname{Proj}$
construction is the algebraic-geometry presentation of
$H^0(\mathbb P^1_K, \mathcal O(1))$. The remaining gap is now
the scheme-theoretic identification
$H^0(\mathrm{Proj}\,K[X, Y], \mathcal O(1)) \simeq K[X, Y]_1$,
a general algebraic-geometry theorem independent of TFPT.

v10 also closes the Stage D backward direction
(\texttt{YukawaStageDExistence.lean}). Given any basis
$b : \texttt{Basis}\,(\texttt{Fin}\,3)\,K\,E$, the determinant
form $b.\texttt{det}$ transports via
\texttt{exteriorPower.alternatingMapLinearEquiv} to an
alternating trilinear form $\omega : \Lambda^3 E \to_{\ell K} K$.
The contraction $C_\omega$ is shown injective by basis-level
matrix-determinant computation and surjective by
$\dim E = \dim(\Lambda^2 E)^* = 3$ and the standard
finite-dimensional injective-iff-surjective theorem. From this
the Stage D characterisation
$\exists\,\omega\;(\textrm{CI}\,\wedge\,\textrm{CS})
  \iff \dim E = 3$ now holds (under $\texttt{Nontrivial}\,E$).

The bundled main theorem packages the three load-bearing
conclusions
\[
(q_-,q_+)=(-2,3), \qquad \tr Y=0, \qquad 6\,Y^2 - Y - \mathbf 1 = \mathbf 0
\]
into a single Lean statement \texttt{hypercharge\_carrier\_packet}
that consumes only structured hypotheses --- with the Yukawa
input now refinable all the way to Stage D, and the boundary
input now arriving via an explicit Calder\'on certificate.

The redesign is governed by one editorial discipline: \emph{the
language must not outrun the formal status}. Two examples:

\begin{itemize}[leftmargin=1.6em, itemsep=2pt]
\item This note never claims to formalise the \emph{Standard-
    Model hypercharge spectrum}, only the $3+2$ carrier
    multiplet. Duals, exterior powers, tensor products --- the
    constructions that lift the carrier to all SM chiral
    multiplets --- are out of scope (\S\ref{sec:limits}).
\item The Yukawa Stage-A certificate is named
    \texttt{YukawaRankCertificate}, not
    \texttt{PrimitiveYukawaTypeCertificate}. The certificate is
    a \emph{rank witness}; it is not a formalisation of
    "primitive indecomposable Yukawa type", which is the proper
    open problem.
\end{itemize}

% ---------------------------------------------------------------- %
\section{What Was Built}
\label{sec:built}
% ---------------------------------------------------------------- %

Twenty-seven Lean modules under
\texttt{experiments/lean4-carrier-rigidity/TfptCarrier/}:

\begin{center}
\footnotesize
\begin{tabularx}{0.98\textwidth}{@{}p{0.28\textwidth}p{0.08\textwidth}Y@{}}
\toprule
\textbf{Module} & \textbf{Lines} & \textbf{Content} \\
\midrule
\texttt{Polarization.lean}        & $\sim150$ & Orthogonal-idempotent algebra; $\mathrm{sixY}^2-\mathrm{sixY}-6=0$. \\
\texttt{InvolutionProjectors.lean}& $\sim200$ & $\varepsilon^2=\mathbf 1 \Rightarrow P_-, P_+$ complete orthogonal idempotents. \\
\texttt{CalderonInterface.lean}   & $\sim110$ & Paper~1 \texttt{CalderonCertificate}: explicit interface wrapper feeding \texttt{CarrierInvolution}. \\
\texttt{CalderonProjector.lean}   & $\sim125$ & Structural derivation: idempotent $\pi^2=\pi$ $\Rightarrow$ $\varepsilon := 2\pi-1$ satisfies $\varepsilon^2 = 1$. \\
\texttt{BoundaryPolarization.lean}& $\sim200$ & v10 upstream: $H = H_+ \oplus H_-$ produces an idempotent endomorphism on $H$ via Mathlib's \texttt{IsCompl.projection}; feeds \texttt{CalderonProjector}. \\
\texttt{MathlibBridge.lean}       & $\sim55$  & \texttt{Polarization} $\leftrightarrow$ \texttt{CompleteOrthogonalIdempotents}. \\
\texttt{LatticeRigidityGeneral.lean}&$\sim200$& General $(m,n)$ primitive rigidity; SM corollary $(3,2)\to(-2,3)$. \\
\texttt{Rigidity.lean}            & $\sim40$  & SM-case wrapper. \\
\texttt{OrientedDeterminantCarrier.lean} & $\sim100$ & \texttt{PrimitiveOrientedDeterminantCarrier (m n)}: trace-zero + orientation + primitivity bundled. \\
\texttt{TraceProjection.lean}     & $\sim100$ & $\tr(aP_-+bP_+) = a\dim E_- + b\dim E_+$; trace-zero corollary. \\
\texttt{DeterminantCharacter.lean}& $\sim130$ & $\det T(\lambda) = \lambda^{m q_-+n q_+}$; universal-$\lambda$ and at-two headlines. \\
\texttt{HiggsIndexShadow.lean}    & $\sim125$ & Stage-A \texttt{HiggsIndexCertificate}; algebraic shadow has dim 2. \\
\texttt{HiggsTopForm.lean}        & $\sim210$ & Stage-B \texttt{HiggsTopForm}: $E_+ \simeq$ space of $\alpha X + \beta Y$ (degree-1 in two formal vars), $\dim = 2$ derived. \\
\texttt{HiggsSchemeCohomologyShadow.lean} & $\sim200$ & v10 upstream: \texttt{DegreeOneTwoVarForm}\,$K$ $\simeq_{\ell K}$ \texttt{MvPolynomial.homogeneousSubmodule}\,(\texttt{Fin}\,2)\,$K$\,1 ($= K[X_0, X_1]_1$, Proj-side $H^0$). \\
\texttt{YukawaRank.lean}          & $\sim135$ & Stage-A \texttt{YukawaRankCertificate} (rank witness, $E_- \simeq K^3$). \\
\texttt{YukawaTopForm.lean}       & $\sim195$ & Stage-B \texttt{YukawaTopForm} ($\Lambda^3 E_- \simeq K$) + rank-3 theorem + characterisation. \\
\texttt{YukawaPrimitive.lean}     & $\sim215$ & Stage-C \texttt{PrimitiveIndecomposableYukawaCoupling} ($E_- \simeq (\Lambda^2 E_-)^*$) + rank-3 theorem + characterisation. \\
\texttt{YukawaTrilinearForm.lean} & $\sim210$ & Stage-D \texttt{YukawaTrilinearForm} ($\omega : \Lambda^3 E_- \to K$) + operational \texttt{ContractionInjective}/\texttt{ContractionSurjective} predicates + constructed contraction + rank-3 theorem (under \texttt{Nontrivial}\,$E$) + unconditional $0\lor3$ transparency theorem. \\
\texttt{YukawaStageDExistence.lean}& $\sim280$ & v10 Stage D backward: $\dim E = 3 \Rightarrow \exists\,\omega$ with CI $\wedge$ CS via $b.\texttt{det}$ + direct $3\times 3$ matrix-determinant computation. \\
\texttt{BoundaryYukawaKernelInterface.lean}& $\sim120$ & v11 abstract typed target: \texttt{BoundaryYukawaKernel} packages $(\omega, \textrm{CI}, \textrm{CS}, \texttt{Nontrivial})$ with characterisation $\iff \dim E_- = 3$. \\
\texttt{SeamWindingInterface.lean}& $\sim130$ & v11 abstract typed target: \texttt{SeamWindingData}\,$(m, n)$ packages the seam-class consequences (trace-zero, orientation, primitivity); bidirectional bridge to \texttt{PrimitiveOrientedDeterminantCarrier}. \\
\texttt{CarrierData.lean}         & $\sim200$ & Bundled main theorems: pair, trace, polynomial, packet. \\
\texttt{Hypercharge.lean}         & $\sim95$  & Concrete $5\times5$ model: $\tr Y=0$ and $6Y^2-Y-\mathbf 1 = \mathbf 0$. \\
\texttt{Sanity.lean}              & $30$      & \texttt{\#eval} smoke tests. \\
\texttt{AxiomCheck.lean}          & $80$      & \texttt{\#print axioms} on every headline theorem. \\
\texttt{AuditCheck.lean}          & $80$      & \texttt{\#check} elaboration of every headline name. \\
\texttt{AuditContract.lean}       & $\sim180$ & Signature-lock \texttt{example} patterns for every headline theorem (strict signature audit). \\
\bottomrule
\end{tabularx}
\end{center}

Plus:
\begin{itemize}[leftmargin=1.5em, itemsep=1pt]
\item \texttt{scripts/audit.sh} --- seven-check hard CI gate
      (build / \texttt{sorry} / \texttt{axiom} /
      \texttt{unsafe-opaque-partial} / axiom-print /
      headline-\texttt{\#check} / signature-lock \texttt{example}).
\item \texttt{lakefile.lean} pinning Lean~4.29.1 and Mathlib~v4.29.1.
\end{itemize}

% ---------------------------------------------------------------- %
\section{Dependency Graph}
\label{sec:graph}
% ---------------------------------------------------------------- %

Every node is a separately formalised Lean theorem; every arrow
is a discharged proof obligation in \texttt{TfptCarrier}. Blue
nodes are unconditional Lean theorems. Orange nodes are Stage-A
structured hypotheses. Green nodes are the bundled conclusions.

\begin{center}
\begin{tikzpicture}[
    node distance=8mm,
    box/.style={rectangle, rounded corners=2pt, draw=tfptblue!65, fill=tfptblue!5,
        align=center, text width=4.0cm, inner sep=4pt, font=\footnotesize\sffamily},
    cert/.style={rectangle, rounded corners=2pt, draw=tfptorange!70!black, fill=tfptorange!8,
        align=center, text width=4.0cm, inner sep=4pt, font=\footnotesize\sffamily},
    bundle/.style={rectangle, rounded corners=2pt, draw=tfptgreen!70!black, fill=tfptgreen!7,
        align=center, text width=4.8cm, inner sep=4pt, font=\footnotesize\sffamily},
    arr/.style={-{Latex[length=2mm]}, thick, draw=tfptgray!85}
]
\node[box] (bndpol) {BoundaryPolarization (v10)\\$H = H_+ \oplus H_-$};
\node[box, below=of bndpol] (calproj) {CalderonProjector (Paper 1, struct.)\\$\pi^2 = \pi$, $\varepsilon := 2\pi-1$};
\node[cert, below=of calproj] (calderon) {CalderonCertificate (Paper 1, iface)\\$\varepsilon^2 = \mathbf 1$};
\node[box, below=of calderon] (eps)    {Boundary involution\\$\varepsilon^2 = \mathbf 1$};
\node[box, below=of eps] (proj)   {InvolutionProjectors\\$\Rightarrow P_-, P_+$ idempotent};
\node[box, right=22mm of calproj] (higgstop) {HiggsTopForm (Stage B Higgs)\\$E_+ \simeq \alpha X{+}\beta Y$};
\node[cert, below=of higgstop] (higgs) {HiggsIndexCertificate (Stage A)\\$\Rightarrow n = \dim E_+ = 2$};
\node[box, below=of higgs]        (trilinear){YukawaTrilinearForm (Stage D)\\$\omega : \Lambda^3 E_- \to K$, NDeg+Prim};
\node[box, below=of trilinear]    (primitive){YukawaPrimitive (Stage C)\\$E_- \simeq (\Lambda^2 E_-)^*$};
\node[box, below=of primitive]    (topform){YukawaTopForm (Stage B)\\$\Lambda^3 E_- \simeq K$};
\node[cert, below=of topform]     (yukawa) {YukawaRankCertificate (Stage A)\\$\Rightarrow m = \dim E_- = 3$};
\node[box, below=of proj]         (trc)    {TraceProjection\\$\tr(aP_-+bP_+) = am+bn$};
\node[box, below=of trc]          (det)    {DeterminantCharacter\\$\det T(\lambda)=1 \Rightarrow m q_-+n q_+=0$};
\node[box, below=of det]          (lat)    {LatticeRigidityGeneral\\$+\,$primitive orient.\,$\Rightarrow (q_-,q_+)$};
\node[bundle, right=12mm of lat]  (carrier){CarrierData\\$(-2,3),\ \tr Y=0,\ 6Y^2{-}Y{-}\mathbf 1=\mathbf 0$};
\node[bundle, below=10mm of carrier] (hyp){Hypercharge ($5\times 5$ model)\\$Y = \mathrm{diag}(-\tfrac13,-\tfrac13,-\tfrac13,\tfrac12,\tfrac12)$};

% Arrows: only the logically correct ones.
\draw[arr] (bndpol) -- (calproj);      % v10: polarisation → idempotent
\draw[arr] (calproj) -- (calderon);    % Paper 1 structural → interface
\draw[arr] (calderon) -- (eps);
\draw[arr] (eps) -- (proj);
\draw[arr] (proj) -- (trc);
\draw[arr] (higgstop) -- (higgs);      % Higgs Stage B → Stage A
\draw[arr] (trilinear) -- (primitive); % Stage D → Stage C
\draw[arr] (primitive) -- (topform);   % Stage C → Stage B
\draw[arr] (topform) -- (yukawa);      % Stage B → Stage A
\draw[arr] (higgs.south west) to[bend right=8] (trc.east);
\draw[arr] (yukawa.west) to[bend right=8] (trc.east);
\draw[arr] (trc) -- (carrier);
\draw[arr] (det) -- (lat);
\draw[arr] (lat) -- (carrier);
\draw[arr] (proj.east) to[bend right=10] (carrier.west);
\draw[arr] (carrier) -- (hyp);
\end{tikzpicture}
\end{center}

\noindent
Three things to note about the graph that the previous version
got wrong, and one new feature:

\begin{itemize}[leftmargin=1.6em, itemsep=2pt]
\item \textbf{No \texttt{InvolutionProjectors} $\to$
    \texttt{LatticeRigidityGeneral} arrow.} The projector
    construction does not imply lattice rigidity; the two
    branches are independent until they meet at
    \texttt{CarrierData}.
\item \textbf{The arrow runs \texttt{DeterminantCharacter} $\to$
    \texttt{LatticeRigidityGeneral}, not the reverse.} The
    determinant character produces the trace-zero equation
    $m q_- + n q_+ = 0$; the rigidity theorem then upgrades it to
    a unique pair when combined with the primitivity and sign
    conditions.
\item \textbf{No \texttt{HiggsIndexCertificate} $\to$
    \texttt{YukawaRankCertificate} arrow.} The two certificates
    discharge different ranks via independent mathematical
    arguments. Both feed into \texttt{CarrierData} separately.
\item \textbf{New in v4: \texttt{YukawaTopForm} $\to$
    \texttt{YukawaRankCertificate} arrow.} The Yukawa input was
    split into two layers. Stage B (\texttt{YukawaTopForm})
    is the algebraic statement $\Lambda^3 E_- \simeq K$, which is
    the precise Lean reading of TFPT Paper~2's "$\Lambda^3 E_- =
    \det E_-$". The Stage A rank witness drops out as a corollary
    via an arbitrary basis.
\item \textbf{New in v5: \texttt{YukawaPrimitive} $\to$
    \texttt{YukawaTopForm} arrow.} The Yukawa input is further
    refined into a Stage C
    \texttt{PrimitiveIndecomposableYukawaCoupling}: a
    contraction iso $E_- \simeq (\Lambda^2 E_-)^*$ plus
    non-triviality.
\item \textbf{New in v6: \texttt{YukawaTrilinearForm} (Stage D)
    $\to$ \texttt{YukawaPrimitive} arrow.} Stage D introduces an
    alternating trilinear form $\omega : \Lambda^3 E_- \to K$ as
    primary datum, together with two explicit operational
    \texttt{Prop} predicates (\texttt{ContractionInjective} and
    \texttt{ContractionSurjective} from v7 onwards, originally
    \texttt{Nondegenerate} and \texttt{Primitive} in v6).
    The contraction $C_\omega$ is \emph{constructed} from
    $\omega$ (no longer data), and the iso conclusion of Stage C
    becomes a Lean theorem under the two predicates and
    \texttt{Nontrivial}\,$E_-$. This is the upstream-most layer
    of the Yukawa tower; the rank-3 conclusion is now derived
    from an alternating trilinear form together with two
    structural predicates and non-triviality.
\item \textbf{New in v6: \texttt{CalderonCertificate} $\to$
    \texttt{CarrierInvolution} arrow.} The Paper 1 input is no
    longer silently consumed: a \texttt{CalderonCertificate} is
    a thin wrapper that names the Calder\'on-induced carrier
    involution as a Lean datum, with a forgetful map into the
    \texttt{InvolutionProjectors} chain.
\item \textbf{New in v7: \texttt{AuditContract.lean} signature
    locks.} Each headline theorem is used in an
    \texttt{example : <exact intended type> := <theorem>}
    pattern, detecting silent weakening of theorem statements
    that \texttt{\#check} would miss.
\item \textbf{New in v7:
    \texttt{PrimitiveOrientedDeterminantCarrier (m n)}.}\ The
    three-part TFPT carrier condition (trace-zero, orientation,
    primitivity) is bundled as a single Lean structure, removing
    the semantic-drift risk of an unstructured triple of
    hypotheses.
\item \textbf{New in v8: \texttt{CalderonProjector} $\to$
    \texttt{CalderonCertificate} arrow.}\ The algebraic step
    $\pi^2 = \pi \Rightarrow \varepsilon^2 = \mathbf 1$ (for
    $\varepsilon := 2\pi - 1$) is now a Lean theorem, upgrading
    the Calder\'on layer from documentary interface to a
    structural algebraic derivation.
\item \textbf{New in v8: \texttt{HiggsTopForm} $\to$
    \texttt{HiggsIndexCertificate} arrow.}\ A Higgs Stage B
    structure identifies $E_+$ with the algebraic shadow of
    $H^0(\mathbb P^1, \mathcal O(1))$, namely
    \texttt{DegreeOneTwoVarForm K}, and derives $\dim E_+ = 2$
    from this polynomial-space iso rather than from a bare
    rank witness.
\item \textbf{New in v9: \texttt{finrank\_E\_eq\_zero\_or\_three}.}\
    Unconditional Stage D transparency theorem: without
    \texttt{Nontrivial}\,$E$ the operational predicates force
    only $\dim E = 0 \lor \dim E = 3$, which is the precise
    mathematical reason the rank-3 slogan must carry the
    \texttt{Nontrivial}\,$E$ hypothesis at the use site.
\item \textbf{New in v10: \texttt{BoundaryPolarization} $\to$
    \texttt{CalderonProjector} arrow.}\ Any $H = H_+ \oplus H_-$
    gives an idempotent endomorphism on $H$ via Mathlib's
    \texttt{IsCompl.projection}. The Calder\'on input is now
    derived from a polarisation datum (one step closer to
    Paper 1's analytic source).
\item \textbf{New in v10:
    \texttt{HiggsSchemeCohomologyShadow} $\to$
    \texttt{HiggsTopForm} arrow.}\ The Higgs Stage-B shadow is
    pinned to Mathlib's
    \texttt{MvPolynomial.homogeneousSubmodule (Fin 2) K 1},
    i.e.\ the Proj-side degree-1 piece of $K[X_0, X_1]$.
\item \textbf{New in v10: \texttt{YukawaStageDExistence}.}\
    Stage D backward direction: $\dim E = 3 \Rightarrow
    \exists\,\omega\;(\textrm{CI} \wedge \textrm{CS})$. The
    Stage D characterisation
    $\exists\,\omega \iff \dim E = 3$ now holds under
    \texttt{Nontrivial}\,$E$.
\end{itemize}

% ---------------------------------------------------------------- %
\section{The Headline Theorems}
\label{sec:theorems}
% ---------------------------------------------------------------- %

We list the headline theorems by graph layer. Every name is
elaborated by \texttt{TfptCarrier/AuditCheck.lean}.

\subsection*{Layer 1 --- Polarization (algebraic core)}

\begin{theorem}[\texttt{Polarization.sixY\_carrier\_polynomial}]
Let $A$ be a ring with two orthogonal idempotents $P_-, P_+$
summing to $\mathbf 1$. Then $\mathrm{sixY}:=-2 P_- + 3 P_+$
satisfies
$\mathrm{sixY}^2 - \mathrm{sixY} - 6 = \mathbf 0$.
\end{theorem}

\subsection*{Layer 0+ --- Calderón structural certificate (Paper 1)}

\begin{theorem}[\texttt{CalderonProjector.eps\_sq}]
Let $A$ be a ring with $2$ invertible and let
$\pi \in A$ satisfy $\pi^2 = \pi$. Define
$\varepsilon := 2\pi - 1$. Then $\varepsilon^2 = \mathbf 1$.
\end{theorem}

\begin{theorem}[\texttt{CalderonProjector.toCalderonCertificate}]
Every \texttt{CalderonProjector} produces a
\texttt{CalderonCertificate} (the Paper 1 interface).
The Calder\'on entry point is now derived from a single
algebraic hypothesis --- the existence of the
boundary-spectral idempotent.
\end{theorem}

\subsection*{Layer 2 --- Involution to projectors}

\begin{theorem}[\texttt{InvolutionProjectors.CarrierInvolution.toPolarization}]
Let $\varepsilon$ be an element of a ring $A$ with $\varepsilon^2 = \mathbf 1$ and
$2$ invertible in $A$. The spectral projectors $P_-, P_+$ defined
by $P_\pm := (\mathbf 1\pm\varepsilon)\cdot\widehat 2$ form a
\texttt{Polarization} of $A$. Each of the four polarization
axioms is discharged directly from $\varepsilon^2 = \mathbf 1$.
\end{theorem}

\subsection*{Layer 3 --- General lattice rigidity}

\begin{theorem}[\texttt{LatticeRigidityGeneral.primitive\_trace\_free\_pair\_general}]
Let $m, n$ be positive integers and $q_-, q_+ \in \mathbb Z$ satisfy
\[
m q_- + n q_+ = 0,\qquad q_- < 0 < q_+,\qquad
\gcd(|q_-|, q_+) = 1.
\]
Then
$q_- = -n/\gcd(m,n)$ and $q_+ = m/\gcd(m,n)$.
\end{theorem}

\begin{corollary}[\texttt{Rigidity.unique\_carrier\_pair}]
For $(m, n) = (3, 2)$ the same three hypotheses force
$(q_-, q_+) = (-2, 3)$.
\end{corollary}

\subsection*{Layer 4 --- Structural trace}

\begin{theorem}[\texttt{TraceProjection.trace\_linear\_combination\_of\_idempotents}]
For any two idempotent endomorphisms $P_-, P_+$ of a finite-
dimensional $K$-vector space $M$, and any $a, b \in K$,
\[
\tr_M(a\cdot P_- + b\cdot P_+)
= a \cdot \dim_K(\operatorname{range} P_-)
+ b \cdot \dim_K(\operatorname{range} P_+).
\]
\end{theorem}

\begin{corollary}[\texttt{TraceProjection.trace\_carrier\_Y\_eq\_zero}]
With $a=-\tfrac13, b=\tfrac12$ and the rank discharges
$\dim E_- = 3, \dim E_+ = 2$, the structural identity gives
$\tr Y = 0$ for $Y = -\tfrac13 P_- + \tfrac12 P_+$.
The hypothesis $[\texttt{CharZero K}]$ is genuinely used here.
\end{corollary}

\subsection*{Layer 5 --- Determinant character}

\begin{theorem}[\texttt{DeterminantCharacter.det\_torusMatrix}]
For nonzero $\lambda$ in a field $K$ and integer charges $q_-, q_+$,
\[
\det T(\lambda)
= \lambda^{\,m \cdot q_- \,+\, n \cdot q_+},
\quad
T(\lambda) :=
\mathrm{diag}\bigl(
\underbrace{\lambda^{q_-},\ldots,\lambda^{q_-}}_{m},
\underbrace{\lambda^{q_+},\ldots,\lambda^{q_+}}_{n}
\bigr).
\]
\end{theorem}

\begin{corollary}[\texttt{DeterminantCharacter.trace\_zero\_of\_det\_one\_forall\_rat}]
If $\det T(\lambda) = 1$ for \emph{every} positive rational
$\lambda \neq 1$, then $m q_- + n q_+ = 0$. The single-point form
$\det T(2) = 1 \Rightarrow m q_- + n q_+ = 0$ is recorded
separately as \texttt{trace\_zero\_of\_det\_one\_at\_two}.
\end{corollary}

\begin{remark}
The Diophantine equation $m q_- + n q_+ = 0$ from the determinant
character is necessary but \emph{not sufficient} to fix $(q_-,
q_+)$: in $(m,n)=(3,2)$ it is satisfied by every pair of the form
$(-2 k, 3 k)$. The primitivity ($\gcd = 1$) and sign ($q_- <
0 < q_+$) conditions are what reduce that infinite family to the
unique pair $(-2, 3)$. The combined hypothesis is therefore
phrased throughout as the \emph{primitive oriented determinant-
preserving} carrier condition.
\end{remark}

\subsection*{Layer 6 --- Stage-A and Stage-B rank certificates}

\begin{theorem}[\texttt{HiggsTopForm.DegreeOneTwoVarForm.finrank\_eq\_two}]
The space of homogeneous degree-1 polynomials in two formal
variables, defined concretely as
$\texttt{DegreeOneTwoVarForm K} := \{(\alpha, \beta) : K \times K\}$
with the natural $K$-module structure, has $K$-rank $2$.
\end{theorem}

\begin{theorem}[\texttt{HiggsTopForm.HiggsTopForm.finrank\_Eplus\_eq\_two}]
A Higgs Stage-B certificate (a linear equivalence
$E_+ \simeq \texttt{DegreeOneTwoVarForm}\,K$) forces
$\dim_K E_+ = 2$. The Stage A → Stage B refinement is a real
algebraic statement: $E_+$ is identified with the polynomial
space, not just with $K^2$.
\end{theorem}

\subsection*{Layer 6 --- Stage-A rank certificates}

\begin{theorem}[\texttt{HiggsIndexShadow.HiggsIndexCertificate.finrank\_Eplus\_eq\_two}]
A \texttt{HiggsIndexCertificate}, defined as a linear equivalence
$E_+ \simeq H^0(\mathbb P^1, \mathcal O(1))^{\mathrm{alg}}
\simeq K^2$, forces $\dim_K E_+ = 2$.
\end{theorem}

\begin{theorem}[\texttt{YukawaRank.YukawaRankCertificate.finrank\_Eminus\_eq\_three}]
A \texttt{YukawaRankCertificate} (Stage A), defined as a linear
equivalence $E_- \simeq K^3$, forces $\dim_K E_- = 3$.
\end{theorem}

\subsection*{Layer 6b --- Yukawa Stage B (algebraic content)}

\begin{theorem}[\texttt{YukawaTopForm.finrank\_lambda3\_eq\_one}]
A \texttt{YukawaTopForm K E}, defined as a linear equivalence
$\Lambda^3 E \simeq K$, forces $\dim_K (\Lambda^3 E) = 1$.
\end{theorem}

\begin{theorem}[\texttt{YukawaTopForm.finrank\_E\_eq\_three}]
For a finite-dimensional $K$-vector space, a \texttt{YukawaTopForm}
forces $\dim_K E = 3$. The proof is the composition
\[
\Lambda^3 E \simeq K
\;\Rightarrow\; \dim \Lambda^3 E = 1
\;\xRightarrow{\text{\texttt{finrank\_eq}}}\; \binom{\dim E}{3} = 1
\;\xRightarrow{\text{\texttt{choose\_eq\_one\_iff}}}\; \dim E = 3.
\]
\end{theorem}

\begin{theorem}[\texttt{YukawaTopForm.ofFinrankEqThree}]
Conversely, any finite-dimensional $K$-vector space with $\dim E = 3$
admits a \texttt{YukawaTopForm}: pick any basis and read off the
resulting one-dimensional $\Lambda^3 E$ as $K$.
\end{theorem}

\begin{theorem}[\texttt{YukawaTopForm.nonempty\_iff\_finrank\_eq\_three}]
Characterisation: for a finite-dimensional $K$-vector space,
\[
\text{Nonempty}\,(\texttt{YukawaTopForm}\,K\,E)
\;\iff\;
\dim_K E = 3.
\]
The Stage B hypothesis is therefore \emph{precisely} the rank-3
condition, expressed in exterior-algebra form.
\end{theorem}

\begin{theorem}[\texttt{YukawaTopForm.toYukawaRankCertificate}]
A Stage B \texttt{YukawaTopForm} produces a Stage A
\texttt{YukawaRankCertificate} by extracting $\dim E = 3$ and
choosing an arbitrary basis.
\end{theorem}

\subsection*{Layer 6c --- Yukawa Stage C (primitive indecomposable
coupling)}

\begin{theorem}[\texttt{PrimitiveIndecomposableYukawaCoupling.finrank\_E\_eq\_three}]
A \texttt{PrimitiveIndecomposableYukawaCoupling K E}, defined as
a linear equivalence $E \simeq (\Lambda^2 E)^*$ plus
non-triviality of $E$, forces $\dim_K E = 3$.

\smallskip\noindent
\emph{Proof sketch.}\ The contraction iso transports dimensions
as
\[
\dim E
  = \dim (\Lambda^2 E)^*
  \overset{\text{dual}}{=} \dim \Lambda^2 E
  \overset{\text{exteriorPower}}{=} \binom{\dim E}{2}
  = \tfrac{n(n-1)}{2}.
\]
The integer equation $n = n(n-1)/2$ with $n \ge 1$ has the
unique solution $n = 3$ (a finite case split / \texttt{omega}).
\end{theorem}

\begin{theorem}[\texttt{PrimitiveIndecomposableYukawaCoupling.toYukawaTopForm}]
A Stage C coupling produces a Stage B
\texttt{YukawaTopForm} via the v4 converse
\texttt{YukawaTopForm.ofFinrankEqThree} applied to
\texttt{finrank\_E\_eq\_three}.
\end{theorem}

\begin{theorem}[\texttt{PrimitiveIndecomposableYukawaCoupling.toYukawaRankCertificate}]
A Stage C coupling produces a Stage A \texttt{YukawaRankCertificate}
through the composition C $\Rightarrow$ B $\Rightarrow$ A.
\end{theorem}

\begin{theorem}[\texttt{PrimitiveIndecomposableYukawaCoupling.ofFinrankEqThree}]
Conversely, any finite-dimensional $K$-vector space with
$\dim E = 3$ admits a Stage C coupling: a contraction iso
$E \simeq (\Lambda^2 E)^*$ exists by dimension count, since both
sides have dimension $3$ (via
\texttt{dual\_finrank\_eq} and \texttt{exteriorPower.finrank\_eq}).
\end{theorem}

\begin{theorem}[\texttt{PrimitiveIndecomposableYukawaCoupling.nonempty\_iff\_finrank\_eq\_three}]
For a finite-dimensional $K$-vector space,
$\text{Nonempty}\,(\texttt{PrimitiveIndecomposableYukawaCoupling}\,K\,E) \iff \dim_K E = 3$.
\end{theorem}

\subsection*{Layer 6d --- Yukawa Stage D (primary trilinear form)}

\begin{theorem}[\texttt{YukawaTrilinearForm.contraction}]
For every alternating trilinear form
$\omega : \Lambda^3 E \to K$, the associated contraction map
\[
C_\omega : E \to (\Lambda^2 E)^*,
\qquad
C_\omega(v)(w \wedge x) := \omega(v \wedge w \wedge x),
\]
is constructible as a \texttt{LinearMap}. The Lean construction
uses \texttt{exteriorPower.alternatingMapLinearEquiv} (to
convert between $(\Lambda^n E)^*$ and the space of $K$-valued
alternating maps on $E^n$) and \texttt{AlternatingMap.curryLeft}
(to split off the first argument).
\end{theorem}

\begin{theorem}[\texttt{YukawaTrilinearForm.ContractionInjective} and \texttt{ContractionSurjective}]
Two explicit operational \texttt{Prop} predicates on $\omega$:
\[
\texttt{ContractionInjective}\,\omega \;:=\; \texttt{Injective}\,C_\omega,
\qquad
\texttt{ContractionSurjective}\,\omega \;:=\; \texttt{Surjective}\,C_\omega.
\]
The TFPT-side phrases "non-degenerate" and "primitive / no
spectator" are docstring synonyms for the two operational
predicates respectively. The combined predicate
\texttt{PrimitiveIndecomposableYukawaCondition} bundles both
(plus \texttt{Nontrivial}\,$E$ at the use site).
\end{theorem}

\begin{theorem}[\texttt{YukawaTrilinearForm.contractionEquivOfBoth}]
A trilinear form $\omega$ satisfying both predicates over a
finite-dimensional vector space induces a linear iso
$E \simeq (\Lambda^2 E)^*$ --- the contraction $C_\omega$ itself,
upgraded to a \texttt{LinearEquiv} via
\texttt{LinearEquiv.ofBijective}.
\end{theorem}

\begin{theorem}[\texttt{YukawaTrilinearForm.toPrimitiveIndecomposableYukawaCoupling}]
A trilinear form with \texttt{ContractionInjective} and
\texttt{ContractionSurjective} on a non-trivial finite-dimensional
vector space yields a Stage C
\texttt{PrimitiveIndecomposableYukawaCoupling}. Together with the
Stage C $\Rightarrow$ Stage B $\Rightarrow$ Stage A chain, this means

\[
\boxed{\;
\omega \,:\,\Lambda^3 E \to K
\;\;\text{(alt.)}\;\;\text{plus}\;\;\texttt{CI}, \texttt{CS}, \texttt{Nontrivial}\,E
\;\;\Longrightarrow\;\;
\dim E = 3.
\;}
\]
\end{theorem}

\begin{theorem}[\texttt{YukawaTrilinearForm.finrank\_E\_eq\_zero\_or\_three} (v9)]
\emph{Unconditional transparency form, no \texttt{Nontrivial}
hypothesis.} For any finite-dimensional $K$-vector space $E$,
a trilinear form $\omega$ satisfying \texttt{ContractionInjective}
and \texttt{ContractionSurjective} forces
\[
\dim_K E = 0 \;\;\lor\;\; \dim_K E = 3.
\]
The two predicates are vacuously satisfied when $E = 0$;
the rank-3 conclusion holds precisely when $E$ is non-trivial.
This theorem records the fact that the operational predicates
do \emph{not} themselves rule out the trivial case, which is the
mathematical reason the rank-3 slogan above carries the
\texttt{Nontrivial}\,$E$ hypothesis.
\end{theorem}

\textbf{What Stage D says about TFPT Paper~2.}
The phrase "primitive indecomposable Yukawa type" in TFPT
Paper~2 contains three semi-formal qualifiers:

\begin{itemize}[leftmargin=1.6em, itemsep=0pt]
\item \emph{primitive} --- the trilinear coupling does not
    factor through a sub-trilinear form;
\item \emph{indecomposable} --- the carrier does not split as
    a direct sum that respects the coupling;
\item \emph{no spectator} --- the closure
    $E_- \otimes \Lambda^2 E_- \to \Lambda^3 E_-$ exhausts
    $E_-$.
\end{itemize}

\noindent
\emph{Lean translation in Stage D.}\ The primary datum is the
alternating trilinear form $\omega$ on $E_-$. The three
qualifiers translate to two explicit operational \texttt{Prop}
predicates on $\omega$:
\texttt{ContractionInjective} (the contraction is injective;
equivalently, no nonzero $v$ has $C_\omega(v) = 0$ ---
\emph{non-degenerate} in TFPT-side prose) and
\texttt{ContractionSurjective} (the contraction is surjective;
equivalently, every linear functional on $\Lambda^2 E_-$ is the
contraction of $\omega$ with some vector --- \emph{primitive /
no spectator} in TFPT-side prose). Together with
\texttt{Nontrivial}\,$E_-$ they imply the Stage C contraction
iso and hence $\dim E_- = 3$.

\textbf{Where the standing Yukawa gap now lives.}
The Lean chain takes the trilinear form $\omega$ together with
the two predicates and \texttt{Nontrivial}\,$E_-$ as input.
The remaining open question is \emph{purely upstream}: does the
TFPT boundary kernel
(\texttt{CalderonCertificate} \emph{plus} representation-
theoretic data that the Yukawa coupling exists) produce a
trilinear form satisfying \texttt{ContractionInjective} and
\texttt{ContractionSurjective}? That question is TFPT theory
work, not Lean engineering. The Lean chain is now flat: there is
no hidden assumption between Stage D and $\dim E_- = 3$.

\subsection*{Layer 7 --- Mathlib bridge}

\begin{theorem}[\texttt{MathlibBridge.Polarization.toCompleteOrthogonalIdempotents}]
A \texttt{TFPT.Carrier.Polarization} on a ring $A$ is canonically
a Mathlib \texttt{CompleteOrthogonalIdempotents} on the pair
$(P_-, P_+)$. A converse map
\texttt{ofCompleteOrthogonalIdempotents} is also provided.
\end{theorem}

\subsection*{Layer 8 --- Bundled main theorems}

The carrier-rigidity payload of the project is the following set
of bundled theorems on
\texttt{CarrierPremises K M}:

\begin{theorem}[\texttt{CarrierPremises.primitive\_pair\_eq\_sm}]
Under the primitive oriented determinant-preserving carrier
condition with rank data $(m, n) = (3, 2)$, $(q_-, q_+) = (-2, 3)$.
\end{theorem}

\begin{theorem}[\texttt{CarrierPremises.trace\_Y\_eq\_zero}]
$\tr_M Y = 0$ for $Y = -\tfrac13 P_- + \tfrac12 P_+$.
\end{theorem}

\begin{theorem}[\texttt{CarrierPremises.carrier\_polynomial\_Y}]
$6\cdot(Y\cdot Y) - Y - \mathbf 1 = \mathbf 0$ in $\mathrm{End}_K M$.
The hypothesis $[\texttt{CharZero K}]$ is used here to provide
$\widehat 6 \in K$.
\end{theorem}

\begin{theorem}[\texttt{CarrierPremises.hypercharge\_carrier\_packet}]
Under \texttt{CarrierPremises} and the primitive oriented
determinant-preserving carrier condition, the three bundled
conclusions above hold \emph{simultaneously}. This is the
top-level claim of the project.
\end{theorem}

\subsection*{Concrete realisation}

\begin{theorem}[\texttt{Hypercharge.trace\_Y, Hypercharge.Y\_carrier\_polynomial}]
On the concrete $5\times 5$ rational diagonal model,
$Y = \mathrm{diag}(-\tfrac13,-\tfrac13,-\tfrac13,\tfrac12,\tfrac12)$
satisfies $\tr Y = 0$ and $6\,Y^2 - Y - \mathbf 1 = \mathbf 0$,
proved by direct per-entry rational arithmetic.
\end{theorem}

% ---------------------------------------------------------------- %
\section{Significance}
\label{sec:significance}
% ---------------------------------------------------------------- %

The two load-bearing chains, with all hypotheses explicitly
visible, are:

\paragraph{Chain~A --- trace identity.}
\[
\text{Calderón certificate}
\;\xRightarrow{\text{Paper 1 interface}}\;
\varepsilon^2 = \mathbf 1
\;\xRightarrow{\text{\texttt{InvolutionProjectors}}}\;
P_-, P_+ \text{ orthogonal idempotents}
\]
together with the rank discharges
\[
\dim E_+ = 2 \quad (\text{\texttt{HiggsIndexCertificate}}),
\qquad
\omega : \Lambda^3 E_- \to K\,\text{+CI+CS+Nontrivial} \quad (\text{Stage D})
\;\Rightarrow\;
\dim E_- = 3
\]
and the determinant-normalised coefficient choice
$Y := -\tfrac13 P_- + \tfrac12 P_+$ give, via
\texttt{TraceProjection},
\[
\tr Y = -\tfrac13 \cdot 3 + \tfrac12 \cdot 2 = 0.
\]
The Yukawa input is now an alternating trilinear form with two
explicit \texttt{Prop} predicates. The contraction iso and the
rank-3 conclusion are both derived theorems.

\paragraph{Chain~B --- integer rigidity.}
The trace-zero condition
\[
m\,q_- + n\,q_+ = 0
\]
follows from the universal-$\lambda$ determinant character. Combined
with the primitivity and orientation conditions
\[
q_- < 0 < q_+,\qquad
\gcd(|q_-|, q_+) = 1,
\]
the general $(m,n)$ rigidity theorem gives
\[
q_- = -\frac{n}{\gcd(m,n)},\qquad q_+ = \frac{m}{\gcd(m,n)}.
\]
For $(m,n) = (3,2)$ this specialises to $(q_-,q_+) = (-2,3)$.
The integer answer is a corollary of a parametric theorem; it is
not hard-coded.

\paragraph{Why both chains matter together.}
The carrier polynomial $6 Y^2 - Y - \mathbf 1 = \mathbf 0$ is a
consequence of the polarization layer alone (it does not need
either rank discharge). The trace identity $\tr Y = 0$ requires
the rank discharges. The integer pair $(-2, 3)$ requires the
primitive oriented determinant-preserving condition together
with $(m, n) = (3, 2)$. The bundled theorem
\texttt{hypercharge\_carrier\_packet} is the unique statement
that requires all upstream pieces simultaneously.

\paragraph{What is \emph{not} claimed.} The $3+2$ multiplet
$Y = \mathrm{diag}(-\tfrac13, -\tfrac13, -\tfrac13, \tfrac12, \tfrac12)$
is the carrier. The \emph{full Standard-Model hypercharge
spectrum} --- all chiral-multiplet hypercharges including duals,
exterior powers, and tensor products --- is not formalised in
this experiment. The title and main theorem both say "$3+2$
hypercharge carrier", not "SM hypercharge spectrum".

% ---------------------------------------------------------------- %
\section{Interface Table: TFPT Paper 2 \texorpdfstring{$\leftrightarrow$}{<->} Lean Status}
\label{sec:interface-table}
% ---------------------------------------------------------------- %

The following table maps each load-bearing claim of TFPT
Paper~2 to its current Lean status and the still-missing
upstream proof. It is the primary defence against scope-creep:
the present note formalises columns~1 and~2; column~3 is
explicitly out of scope.

\smallskip
\begin{center}
\footnotesize
\begin{tabularx}{\textwidth}{@{}p{0.30\textwidth}p{0.30\textwidth}Y@{}}
\toprule
\textbf{TFPT Paper 2 claim} & \textbf{Lean status (v11)} & \textbf{Missing upstream proof} \\
\midrule
$\varepsilon^2 = \mathbf 1$ on $E$ &
  Two structural layers: \texttt{CalderonProjector}
  (idempotent $\pi$ gives $\varepsilon := 2\pi-1$ with $\varepsilon^2 = 1$)
  and \texttt{BoundaryPolarization} (any
  $H = H_+ \oplus H_-$ gives such an idempotent via
  \texttt{IsCompl.projection}); feeds
  \texttt{CalderonCertificate}. &
  Analytic construction of a boundary spectral decomposition
  $H = H_+ \oplus H_-$ from Paper 1's elliptic Calder\'on
  operator on Sobolev spaces. \\[3pt]
$\dim E_+ = 2$ &
  Stages A, B, and B+ all formalised
  (\texttt{HiggsIndexCertificate},
  \texttt{HiggsTopForm},
  \texttt{HiggsSchemeCohomologyShadow});
  \texttt{DegreeOneTwoVarForm}\,$K$ $\simeq_{\ell K}$
  \texttt{MvPolynomial.homogeneousSubmodule (Fin 2) K 1}
  $= K[X_0, X_1]_1$, the algebraic-geometry "Proj-side"
  presentation of $H^0(\mathbb P^1, \mathcal O(1))$ &
  Scheme-theoretic identification
  $H^0(\mathrm{Proj}\,K[X,Y], \mathcal O(1)) \simeq K[X, Y]_1$
  (a general algebraic-geometry theorem, independent of TFPT). \\[3pt]
$\dim E_- = 3$ &
  Stage A--D forward + Stage D backward
  (\texttt{YukawaStageDExistence}, v10):
  $\dim E_- = 3 \Rightarrow \exists\,\omega$
  with \texttt{ContractionInjective} and
  \texttt{ContractionSurjective}, via $b.\texttt{det}$ and
  direct $3{\times}3$ matrix computation. The Stage D
  characterisation
  $\exists\,\omega\;(\textrm{CI} \wedge \textrm{CS})
    \iff \dim E_- = 3$ holds under \texttt{Nontrivial}\,$E_-$. &
  TFPT-internal construction of $\omega$ from boundary kernel
  and representation data (the existence is now algebraically
  closed by basis-determinant; what remains is whether TFPT's
  boundary derivation produces \emph{this particular}
  $\omega$). \\[3pt]
$3 q_- + 2 q_+ = 0$ &
  Theorem from determinant character of carrier torus
  (\texttt{DeterminantCharacter}) &
  Seam-winding $[u_\Sigma] = 1$ produces the determinant-
  preserving condition. \\[3pt]
$(q_-, q_+) = (-2, 3)$ &
  Theorem from general $(m,n)$ lattice rigidity
  (\texttt{LatticeRigidityGeneral}) with
  \texttt{PrimitiveOrientedDeterminantCarrier} as input &
  Orientation $q_- < 0 < q_+$ and primitivity
  $\gcd = 1$ from seam-class data. \\[3pt]
$Y = \mathrm{diag}(-\tfrac13,-\tfrac13,-\tfrac13,\tfrac12,\tfrac12)$ &
  Bundled theorem
  (\texttt{CarrierPremises.hypercharge\_carrier\_packet}) plus
  concrete $5\times 5$ model (\texttt{Hypercharge}) &
  None (this is the conclusion). \\[3pt]
$\tr_E Y = 0$ &
  Bundled theorem
  (\texttt{CarrierPremises.trace\_Y\_eq\_zero}) &
  None. \\[3pt]
$6\, Y^2 - Y - \mathbf 1 = \mathbf 0$ &
  Bundled theorem
  (\texttt{CarrierPremises.carrier\_polynomial\_Y}) in
  $\mathrm{End}_K M$ &
  None. \\[3pt]
$S^+ = \Lambda^{\mathrm{even}} E$, family $\nu^c$ &
  \emph{Not formalised} &
  Exterior-power packet construction; chiral-family decomposition. \\[3pt]
$N_{\mathrm{fam}} = 3$ &
  \emph{Not formalised} &
  Seam winding $\to$ family count. \\[3pt]
$G_{\mathrm{phys}} = (SU(3)\times SU(2)\times U(1)_Y)/\mathbb Z_6$ &
  \emph{Not formalised} &
  Stabilizer theorem for the carrier datum; $\mathbb Z_6$
  quotient theorem. \\[3pt]
SM hypercharges on all chiral multiplets &
  \emph{Not formalised} &
  Derived representations: dual, exterior powers,
  tensor products of the carrier. \\[3pt]
$b_1 = 41/10$, $\Omega_{\mathrm{adm}} = 48$, $N_\Phi = 1$ &
  \emph{Not formalised} &
  Index ledger; representation audit. \\
\bottomrule
\end{tabularx}
\end{center}

\smallskip\noindent
The pattern is clear. The Lean note formalises the
\emph{carrier-rigidity algebra} of Paper 2 in full; the
\emph{TFPT-internal derivations} feeding that algebra remain
upstream work; and the \emph{downstream SM-packet content}
(derived representations, $\mathbb Z_6$ quotient, family
counting, index ledger) is intentionally out of scope.

% ---------------------------------------------------------------- %
\section{Limitations}
\label{sec:limits}
% ---------------------------------------------------------------- %

\paragraph{Yukawa: Stage D resolves the v5 critique.}
The v5 critique of Stage C was that it bundled the contraction
iso as data, so it was equivalent (in finite dimensions) to
$\dim E = 3$ and did not "explain" the rank. v6 resolved this
by introducing Stage D; v7 renamed its operational predicates
to keep semantic-load words out of the Lean interface:

\[
\underbrace{\omega : \Lambda^3 E_- \to K}_{\text{primary datum}}
\;\;
\underbrace{+\;\texttt{ContractionInjective}\;+\;\texttt{ContractionSurjective}\;+\;\texttt{Nontrivial}\,E_-}_{\text{operational predicates and non-triviality}}
\;\xRightarrow{\text{Lean theorem}}\;
\dim E_- = 3.
\]

The contraction $C_\omega$ is a \emph{constructed} linear map;
the iso conclusion is a Lean theorem under the two predicates
and non-triviality of $E_-$; the rank-3 statement follows from
the iso via elementary dimension counting. There is no longer
any structural ambiguity at the Yukawa interface in Lean.

\paragraph{Why \texttt{Nontrivial}\,$E_-$ is necessary.}
Without the non-triviality hypothesis the rank-3 slogan would be
false: for $E_- = 0$ the spaces $\Lambda^2 E_-$ and $\Lambda^3
E_-$ are both zero, the contraction $C_\omega : 0 \to 0$ is
trivially bijective, both operational predicates hold vacuously,
but $\dim E_- = 0 \ne 3$. The transparency theorem
\texttt{finrank\_E\_eq\_zero\_or\_three} (v9) records the
unconditional disjunction $\dim E_- = 0 \lor \dim E_- = 3$;
the rank-3 conclusion is the case under
\texttt{Nontrivial}\,$E_-$. Both signatures are signature-locked
in \texttt{AuditContract.lean}.

\paragraph{Stage D characterisation, both directions formalised (v10).}
For all four Yukawa stages the note now proves the full
characterisation
$\text{Nonempty}\,(\text{Stage}) \iff \dim E = 3$.
For Stage D, under \texttt{Nontrivial}\,$E$:
\[
\text{Nonempty}\,\bigl\{\,Y : \texttt{YukawaTrilinearForm}\,K\,E
  \;\big|\; Y.\texttt{ContractionInjective}
  \land Y.\texttt{ContractionSurjective}\,\bigr\}
\;\iff\;
\dim E = 3.
\]
The forward direction is
\texttt{finrank\_E\_eq\_three\_of\_condition}; the backward
direction is \texttt{YukawaStageDExistence.exists\_yukawaTrilinearForm\_of\_finrank\_eq\_three},
proved via \texttt{Module.Basis.det} transported through
\texttt{exteriorPower.alternatingMapLinearEquiv} and direct
$3\times 3$ matrix-determinant computation. The combined
characterisation is \texttt{stage\_D\_iff\_finrank\_eq\_three}.

\paragraph{What remains is upstream TFPT theory work.}
The Lean chain takes the Stage D quadruple
$(\omega, \texttt{ContractionInjective},
\texttt{ContractionSurjective}, \texttt{Nontrivial}\,E_-)$
as input. The remaining open question is whether TFPT itself
produces such a quadruple from the boundary kernel. Concretely:

\begin{itemize}[leftmargin=1.6em, itemsep=1pt]
\item \emph{Existence of $\omega$.}\ Paper 2's Yukawa-type
    argument should produce an explicit alternating trilinear
    form on $E_-$. This is a representation-theoretic statement
    about the boundary kernel output, not a Lean engineering
    statement.
\item \emph{Verification of the two predicates.}\ Once $\omega$
    is constructed, the two predicates need to be checked.
    These checks are local algebraic statements (injectivity:
    no nonzero $v$ vanishes against all wedges; surjectivity:
    every dual functional on $\Lambda^2 E_-$ is realised as a
    contraction).
\end{itemize}

Neither task requires modifying the present Lean chain. They are
TFPT mathematical work, properly upstream of this experiment.

\paragraph{Remaining limitations of the Lean chain itself.}
\begin{itemize}[leftmargin=1.6em, itemsep=1pt]
\item \textbf{Calder\'on: two structural layers formalised,
    analytic input of the polarisation open.}\
    \texttt{CalderonProjector} proves $\pi^2 = \pi
    \Rightarrow \varepsilon^2 = \mathbf 1$, and
    \texttt{BoundaryPolarization} (v10) shows that any
    direct-sum decomposition $H = H_+ \oplus H_-$ of a
    $K$-module produces such an idempotent via Mathlib's
    \texttt{IsCompl.projection}. What remains is the analytic
    Paper 1 theorem producing the boundary spectral
    decomposition itself --- a question of boundary elliptic
    theory on Sobolev spaces, independent of the carrier
    algebra.
\item \textbf{Higgs: algebraic Stage B and scheme-cohomology
    shadow formalised, Proj construction open.}\
    \texttt{HiggsTopForm} identifies $E_+$ with the algebraic
    shadow of $H^0(\mathbb P^1, \mathcal O(1))$ (degree-1 binary
    forms), and v10's \texttt{HiggsSchemeCohomologyShadow}
    refines this further: the algebraic shadow is now exhibited
    as a Mathlib-native object,
    \texttt{MvPolynomial.homogeneousSubmodule (Fin 2) K 1}.
    What remains is the general algebraic-geometry theorem
    $H^0(\mathrm{Proj}\,K[X,Y], \mathcal O(1)) \simeq K[X, Y]_1$,
    which lives in \texttt{Mathlib.AlgebraicGeometry.*} and is
    independent of TFPT.
\item \textbf{Yukawa: Stage D both directions formalised.}\
    v10's \texttt{YukawaStageDExistence} closes the backward
    direction:
    $\dim E_- = 3 \Rightarrow \exists\,\omega$ with
    \texttt{ContractionInjective} and
    \texttt{ContractionSurjective}, via $b.\texttt{det}$ and
    direct $3{\times}3$ matrix-determinant computation. The
    characterisation
    $\exists\,\omega\;(\textrm{CI} \wedge \textrm{CS})
      \iff \dim E_- = 3$ holds under
    $\texttt{Nontrivial}\,E_-$. What remains is the
    TFPT-internal question: does the boundary kernel produce
    \emph{the same} $\omega$ (or one equivalent up to action of
    the carrier symmetries)?
\item \textbf{Scope discipline.}\ The Lean chain formalises the
    $3+2$ carrier hypercharge. It does \emph{not} formalise the
    full SM hypercharge spectrum on duals, exterior powers, and
    tensor products. The $\mathbb Z_6$ quotient and the
    derived-representation extensions are out of scope. See
    \S\ref{sec:next}.
\end{itemize}

\paragraph{Other items deliberately out of scope.}
\begin{itemize}[leftmargin=1.6em, itemsep=0pt]
\item Full $\mathbb Z_6$-quotient theorem for
    $G_{\mathrm{phys}} = (SU(3)\times SU(2)\times U(1)_Y)/\mathbb Z_6$.
\item Derived Standard-Model representations (dual, exterior
    powers, tensor products); the diagonal formalised here is
    the carrier hypercharge, not the full SM hypercharge
    spectrum.
\item Analytic content of Paper~1 (Calder\'on polarisation,
    boundary elliptic theory).
\end{itemize}

% ---------------------------------------------------------------- %
\section{How It Was Verified}
% ---------------------------------------------------------------- %

The audit gate \texttt{scripts/audit.sh} enforces seven hard
checks on every commit:

\begin{enumerate}[leftmargin=1.6em, itemsep=0pt]
\item \texttt{lake build} succeeds (currently $\sim 1\,800$ jobs,
      $\sim 10$~s on top of the Mathlib cache).
\item No \texttt{sorry} or \texttt{admit} in the source.
\item No \texttt{axiom} or \texttt{constant} declarations.
\item No \texttt{unsafe}, \texttt{opaque}, or \texttt{partial}
      declarations.
\item \texttt{\#print axioms} on every headline theorem returns a
      subset of $\{\texttt{propext},\ \texttt{Classical.choice},\
      \texttt{Quot.sound}\}$.
\item Every headline name in
      \texttt{TfptCarrier/AuditCheck.lean} elaborates via
      Lean's \texttt{\#check}. This catches renames, deletions,
      and elaboration failures that a text search would miss.
\item Every headline theorem in
      \texttt{TfptCarrier/AuditContract.lean} elaborates inside
      an
      \texttt{example : <exact intended type> := <theorem>}
      pattern. This is strictly stronger than \texttt{\#check}:
      if a future refactor silently weakens a theorem's
      conclusion without renaming it, the
      \texttt{example} fails to typecheck and the audit fails.
\end{enumerate}

Current verbatim output of \texttt{audit.sh}:

\begin{lstlisting}
=== Build ===
  [PASS] lake build succeeded

=== sorry / admit ===
  [PASS] no sorry or admit found

=== axiom / constant declarations ===
  [PASS] no axiom or constant declarations

=== unsafe / opaque / partial declarations ===
  [PASS] no unsafe / opaque / partial declarations

=== axiom audit (#print axioms) ===
  [PASS] all theorems use only [propext, Classical.choice, Quot.sound]

=== headline theorems (Lean #check) ===
  [PASS] AuditCheck.lean elaborated; all headline #checks succeeded

=== theorem signatures (Lean example signature-lock) ===
  [PASS] AuditContract.lean elaborated; all signature locks hold

=== Final verdict ===
  AUDIT: PASS
\end{lstlisting}

\texttt{\#eval} smoke tests in \texttt{Sanity.lean}:

\begin{lstlisting}
info: TfptCarrier/Sanity.lean:27:0: 0
info: TfptCarrier/Sanity.lean:30:0:
  !![0, 0, 0, 0, 0; 0, 0, 0, 0, 0; 0, 0, 0, 0, 0;
     0, 0, 0, 0, 0; 0, 0, 0, 0, 0]
\end{lstlisting}

% ---------------------------------------------------------------- %
\section{Changelog}
\label{sec:changes}
% ---------------------------------------------------------------- %

\subsection*{v11 --- review-pass closure and typed targets for
the remaining upstream theorems (the present revision)}

Relative to v10, this version is a closure-and-naming pass. It
addresses the v10 reviewer's three actionable findings:

\begin{itemize}[leftmargin=1.6em, itemsep=2pt]
\item \textbf{Editorial consistency.}\ The stale "Stage D
    characterisation as targeted future work" paragraph and the
    matching "Stage D backward" entry in \S\ref{sec:next}, both
    inherited from v9, are removed: v10's
    \texttt{YukawaStageDExistence} closed that gap.

\item \textbf{Wording.}\ "Upstream interfaces structurally
    closed" is reworded to "closed to the next algebraic
    layer", honestly reflecting that Calder\'on, Higgs Stage B,
    and Yukawa Stage D backward are now Lean-proved
    one-layer-up structural derivations, while the analytic
    / TFPT-internal content of \emph{that next layer up}
    remains open.

\item \textbf{Three-part claim contract.}\ The
    \texttt{tfptclaimcontract} environment is restructured to
    separate (a) direct \texttt{CarrierPremises} inputs,
    (b) Lean-proved upstream refinements, (c) open
    TFPT-origin theorems. The third group now lists three
    typed Lean targets for the open content.

\item \textbf{New module
    \texttt{BoundaryPolarization.ofCarrierInvolution}.}\
    The converse of v10's
    \texttt{BoundaryPolarization.toCarrierInvolution}: given
    a carrier involution $\varepsilon^2 = \mathbf 1$ on
    \texttt{Module.End K H}, the +1 and -1 eigenspaces form a
    \texttt{BoundaryPolarization}. This makes the Paper-1
    interface a Lean-proved \emph{algebraic equivalence}: the
    polarisation picture and the involution picture are two
    presentations of the same finite-dimensional carrier
    datum. \emph{Not} the analytic Paper 1 derivation; the
    construction of the polarisation from boundary elliptic
    theory remains upstream.

\item \textbf{New module \texttt{SeamWindingInterface.lean}.}\
    Introduces \texttt{SeamWindingData}\,$(m, n)$, an abstract
    typed target for TFPT Paper 2's seam-class theorem
    "$[u_\Sigma] = 1$ produces a primitive oriented
    determinant-preserving carrier". Provides bidirectional
    bridges to and from \texttt{PrimitiveOrientedDeterminantCarrier}:
    the seam-winding interface is the existing rigidity input,
    renamed to make the upstream TFPT theorem visible.

\item \textbf{New module \texttt{BoundaryYukawaKernelInterface.lean}.}\
    Introduces \texttt{BoundaryYukawaKernel}\,$K\,E$,
    \emph{packaging}\ the Yukawa Stage D quadruple
    $(\omega, \texttt{CI}, \texttt{CS}, \texttt{Nontrivial})$.
    This is \emph{not} a new physical derivation; it is a
    typed renaming of the Stage D condition, supplied so that
    the open TFPT theorem
    "boundary kernel produces a canonical $\omega$" has an
    explicit Lean-side conclusion. Forward direction
    \texttt{finrank\_E\_eq\_three} reuses v10's Stage D
    discharge; backward direction \texttt{ofFinrankEqThree}
    reuses v10's \texttt{YukawaStageDExistence} to give the
    characterisation
    \texttt{nonempty\_iff\_finrank\_eq\_three}. The
    \emph{canonicity} of $\omega$ from the boundary kernel
    --- whether the TFPT-internal construction produces the
    same form as the basis-built witness, up to carrier
    symmetries --- remains the upstream physics question.

\item \textbf{AuditCheck, AxiomCheck, AuditContract extended.}\
    Fifteen new headline names added across the three new
    modules; all depend only on
    $\{\texttt{propext},\ \texttt{Classical.choice},\
    \texttt{Quot.sound}\}$. \texttt{audit.sh} returns
    \texttt{AUDIT:\ PASS}.
\end{itemize}

\smallskip
After v11, the boundaries of the Lean-side carrier-rigidity
chain are sharply visible. The three open TFPT-origin theorems
each have a typed Lean target:

\[
\begin{aligned}
\text{Paper 1 analytic boundary theory}
&\;\Rightarrow\; \texttt{BoundaryPolarization}, \\
\text{TFPT boundary kernel + reps}
&\;\Rightarrow\; \texttt{BoundaryYukawaKernel}, \\
\text{seam class }[u_\Sigma] = 1
&\;\Rightarrow\; \texttt{SeamWindingData}.
\end{aligned}
\]

\subsection*{v10 --- upstream interfaces closed to the next
algebraic layer}

Relative to v9, this version closes the three standing upstream
gaps documented in the v9 limitations section. Each gap is
replaced by a Mathlib-native structural Lean module:

\begin{itemize}[leftmargin=1.6em, itemsep=2pt]
\item \textbf{New module \texttt{BoundaryPolarization.lean}.}\
    Introduces the structure
    \texttt{BoundaryPolarization K H := \{Hplus Hminus :
    Submodule K H, isCompl Hplus Hminus\}}, the algebraic
    skeleton of the Paper 1 polarisation. The projection
    onto $H_+$ along $H_-$ is an idempotent endomorphism
    (via Mathlib's \texttt{Submodule.IsCompl.projection} and
    its \texttt{IsIdempotentElem} lemma); a bridge
    \texttt{toCalderonProjector} feeds the
    \texttt{CalderonProjector} chain in
    \texttt{Module.End K H}, using the standard
    \texttt{Invertible (2 : \texttt{End} K M)} instance from
    Mathlib's \texttt{QuadraticForm.Basic}.

\item \textbf{New module \texttt{HiggsSchemeCohomologyShadow.lean}.}\
    Constructs a \texttt{LinearEquiv} between
    \texttt{DegreeOneTwoVarForm K} (the algebraic shadow used
    by \texttt{HiggsTopForm}) and Mathlib's
    \texttt{MvPolynomial.homogeneousSubmodule (Fin 2) K 1}, the
    degree-1 piece of the bivariate polynomial ring. The latter
    is the algebraic-geometry "Proj-side" presentation of
    $H^0(\mathbb P^1, \mathcal O(1))$. The proof uses
    \texttt{MvPolynomial.homogeneousSubmodule\_one\_eq\_span\_X}
    and \texttt{Submodule.mem\_span\_range\_iff\_exists\_fun}
    to extract coefficients; injectivity is by direct
    extraction of $X_i$-coefficients.

\item \textbf{New module \texttt{YukawaStageDExistence.lean}.}\
    Closes the Stage D backward direction:
    $\dim E = 3 \Rightarrow \exists\,\omega$ with
    \texttt{ContractionInjective} and
    \texttt{ContractionSurjective}. The construction uses
    Mathlib's \texttt{Module.Basis.det} (the determinant of
    vectors in a basis) transported via
    \texttt{exteriorPower.alternatingMapLinearEquiv} into a
    Yukawa form on $\Lambda^3 E$. Injectivity is proved by
    direct $3\times 3$ matrix-determinant computation in
    the basis; surjectivity follows from
    $\dim E = \dim(\Lambda^2 E)^* = 3$ and Mathlib's
    \texttt{LinearMap.injective\_iff\_surjective\_of\_finrank\_eq\_finrank}.
    The Stage D characterisation
    $\exists\,\omega\;(\textrm{CI}\,\wedge\,\textrm{CS})
      \iff \dim E = 3$ now holds (under
    \texttt{Nontrivial}\,$E$).

\item \textbf{AuditCheck, AxiomCheck, AuditContract extended.}\
    Eighteen new headline names added; all depend only on
    $\{\texttt{propext},\ \texttt{Classical.choice},\
    \texttt{Quot.sound}\}$. \texttt{audit.sh} returns
    \texttt{AUDIT:\ PASS} with seven hard checks.
\end{itemize}

\smallskip
After v10, the conditional Lean theorem is

\[
\boxed{\;
\begin{aligned}
\texttt{BoundaryPolarization}\;(H = H_+ \oplus H_-)
&\;\Rightarrow\; \pi^2 = \pi
\;\Rightarrow\; \varepsilon^2 = \mathbf 1
\;\Rightarrow\; P_\pm \text{ idempotent} \\
K[X_0, X_1]_1 \simeq \texttt{DegreeOneTwoVarForm}
&\;\Rightarrow\; \dim E_+ = 2 \\
\dim E_- = 3\;(\text{with }\texttt{Nontrivial}\,E_-)
&\;\iff\; \exists\,\omega\;(\textrm{CI} \wedge \textrm{CS}) \\
m q_- + n q_+ = 0,\ q_- < 0 < q_+,\ \gcd = 1
&\;\Rightarrow\; (q_-, q_+) = (-n/d, m/d)
\end{aligned}
\;}
\]

with every step a Lean theorem, no \texttt{sorry}, all axioms
$\subseteq\{\texttt{propext},\ \texttt{Classical.choice},\
\texttt{Quot.sound}\}$.

\subsection*{v9 --- review-readiness pass and Stage D transparency}

Relative to v8, this version is a consistency-and-transparency
revision rather than a structural one. It addresses internal
version drift identified in the v8 review and adds one
unconditional transparency theorem at the Stage D interface:

\begin{itemize}[leftmargin=1.6em, itemsep=2pt]
\item \textbf{New theorem
    \texttt{YukawaTrilinearForm.finrank\_E\_eq\_zero\_or\_three}.}\
    For any finite-dimensional $K$-vector space (no
    \texttt{Nontrivial} hypothesis) a trilinear form with
    \texttt{ContractionInjective} and
    \texttt{ContractionSurjective} satisfies
    $\dim E = 0 \lor \dim E = 3$. The two predicates are
    vacuously satisfied at $E = 0$, which is the mathematical
    reason the existing rank-3 theorem requires
    \texttt{Nontrivial}\,$E$.

\item \textbf{Stage D slogans and statements carry
    \texttt{Nontrivial}\,$E$ explicitly.}\ The frontbox,
    abstract, headline boxed equations, and interface table now
    state the rank-3 conclusion under \texttt{Nontrivial}\,$E$.
    The Lean theorem already carried this hypothesis; v9 fixes
    the documentation to match.

\item \textbf{Predicate naming unified throughout the note.}\
    All operational mentions use \texttt{ContractionInjective}
    and \texttt{ContractionSurjective}. The TFPT-side phrases
    "non-degenerate" and "primitive / no spectator" are
    retained only as docstring synonyms in mathematical prose.

\item \textbf{Audit gate uniformly described as seven checks.}\
    The earlier v6 wording "six checks" is removed; the
    \texttt{AuditContract.lean} signature-lock check is
    consistently the seventh.

\item \textbf{Limitations section rewritten to match v8/v9.}\
    The previous "Higgs Stage A only" and "Calder\'on interface
    is documentary only" wording is replaced by the accurate v9
    status: Higgs has a Stage B algebraic shadow, Calder\'on has
    an algebraic idempotent-to-involution layer, and the
    remaining open content is precisely identified.

\item \textbf{Calder\'on and Higgs wording sharpened.}\ The note
    no longer claims that Paper 1 or scheme cohomology is
    formalised. Calder\'on is presented as the algebraic step
    $\pi^2 = \pi \Rightarrow \varepsilon^2 = \mathbf 1$; Higgs is
    presented as the algebraic shadow $E_+ \simeq
    \texttt{DegreeOneTwoVarForm K}$.

\item \textbf{AxiomCheck, AuditCheck, AuditContract updated.}\
    The three audit modules now include the new
    \texttt{finrank\_E\_eq\_zero\_or\_three} headline. All
    audit checks continue to return \texttt{AUDIT:\ PASS}.
\end{itemize}

\subsection*{v8 --- Calderón structural + Higgs Stage B}

Relative to v7, this version upgrades two upstream interfaces
from documentary wrappers to structural derivations, narrowing
the "out of scope" column of the interface table:

\begin{itemize}[leftmargin=1.6em, itemsep=2pt]
\item \textbf{New module \texttt{CalderonProjector.lean}.}\
    Introduces the structure
    \texttt{CalderonProjector A := \{\(\pi : A\), \(\pi^2 = \pi\)\}},
    the algebraic skeleton of the Paper 1 Calder\'on
    construction. Defines the spectral reflection
    $\varepsilon := 2\pi - 1$ and \emph{proves}
    $\varepsilon^2 = \mathbf 1$ via the elementary identity
    $(2\pi - 1)^2 = 4\pi^2 - 4\pi + 1 = 1$. The bridge
    \texttt{toCalderonCertificate} feeds the structural result
    into the existing
    \texttt{CalderonInterface.CalderonCertificate} wrapper, so
    the Paper 1 input is now derived from a single algebraic
    hypothesis: \emph{the existence of an idempotent}.

\item \textbf{New module \texttt{HiggsTopForm.lean}.}\
    Introduces the structure
    \texttt{HiggsTopForm K Eplus :=
            \{Eplus $\simeq$ DegreeOneTwoVarForm K\}},
    where
    \texttt{DegreeOneTwoVarForm K := \{coeffX coeffY : K\}}
    is the algebraic shadow of
    $H^0(\mathbb P^1_K, \mathcal O(1))$ --- the space of
    homogeneous degree-1 polynomials $\alpha X + \beta Y$ in
    two formal variables over $K$. The Stage-B rank theorem
    $\dim E_+ = 2$ follows mechanically (\texttt{finrank\_eq\_two}
    on \texttt{DegreeOneTwoVarForm}, transported via the
    \texttt{HiggsTopForm} equivalence). A bridge
    \texttt{toHiggsIndexCertificate} produces the Stage A
    witness.

\item \textbf{Updated audit gate.}\
    \texttt{AxiomCheck.lean}, \texttt{AuditCheck.lean}, and
    \texttt{AuditContract.lean} extended with five new headline
    theorems. All five depend only on
    $\{\texttt{propext},\ \texttt{Classical.choice},\
    \texttt{Quot.sound}\}$. \texttt{audit.sh} continues to
    return \texttt{AUDIT:\ PASS} with seven hard checks.
\end{itemize}

\smallskip
The Lean-side picture from v8 onwards is:

\[
\begin{aligned}
\text{Calder\'on projector}\; (\pi^2 = \pi)
&\;\Rightarrow\; \varepsilon^2 = \mathbf 1 \;\Rightarrow\; P_-, P_+ \text{ idempotent}\\
\text{Higgs Stage B}\; (E_+ \simeq \alpha X + \beta Y)
&\;\Rightarrow\; \dim E_+ = 2\\
\text{Yukawa Stage D}\; (\omega, \texttt{CI}, \texttt{CS}, \texttt{Nontrivial}\,E_-)
&\;\Rightarrow\; \dim E_- = 3
\end{aligned}
\]

with all three discharges now derived from \emph{structural
hypotheses}, not bare numerical / linear-equivalence
assumptions.

\paragraph{Remaining Lean lemma gap.}
The Stage D characterisation backward direction --- existence
of a $(\omega, \texttt{ContractionInjective},
\texttt{ContractionSurjective})$ triple from $\dim E = 3$ ---
requires showing that the wedge map
$E \to \mathrm{Hom}(\Lambda^2 E, \Lambda^3 E)$, $v \mapsto (\eta \mapsto v \wedge \eta)$,
is injective in dimension 3. The cleanest formalisation route
extends a non-zero vector to a basis and uses
\texttt{exteriorPower.iotaMulti\_family\_linearIndependent\_ofBasis}
to deduce $v \wedge e_1 \wedge e_2 \neq 0$. Mathlib does not
yet have a direct lemma for the non-degeneracy of the wedge
pairing $\Lambda^p E \otimes \Lambda^{n-p} E \to \Lambda^n E$,
so this is a single missing infrastructure lemma; the
remainder of the proof is bookkeeping.

\subsection*{v7 --- Signature lock, predicate rename, interface
discipline}

Relative to v6, this version hardens the audit interface,
neutralises a semantically-loaded naming choice, and tightens
the editorial framing around what the note does and does not
claim:

\begin{itemize}[leftmargin=1.6em, itemsep=2pt]
\item \textbf{New \texttt{AuditContract.lean} module.}\ Each
    headline theorem is now used in an
    \texttt{example : <exact intended type> := <theorem>}
    pattern. This is strictly stronger than the \texttt{\#check}
    audit of \texttt{AuditCheck.lean}: it detects silent
    weakening of theorem statements that \texttt{\#check} would
    miss (since a weakened theorem still elaborates under its
    own name). The \texttt{audit.sh} gate now has \emph{seven}
    checks; the new seventh requires
    \texttt{AuditContract.lean} to elaborate.

\item \textbf{Stage D predicate rename.}\
    \texttt{YukawaTrilinearForm.Nondegenerate} $\to$
    \texttt{ContractionInjective} and
    \texttt{YukawaTrilinearForm.Primitive} $\to$
    \texttt{ContractionSurjective}. The new names are purely
    operational --- they describe the bijectivity of
    $C_\omega$ --- and leave the semantic content
    ("non-degenerate", "primitive") to the docstrings. A
    combined predicate
    \texttt{PrimitiveIndecomposableYukawaCondition} is
    available for the TFPT-side phrase. This removes the
    "primitive renamed to surjective" attack surface.

\item \textbf{New
    \texttt{OrientedDeterminantCarrier.lean} module.}\
    The three-part TFPT carrier condition --- trace-zero,
    orientation, primitivity --- is now packaged as a single
    Lean structure
    \texttt{PrimitiveOrientedDeterminantCarrier (m n : \(\mathbb N\))}.
    Producing such a structure is the formal target for
    a future TFPT-internal theorem
    "seam winding $[u_\Sigma] = 1$ implies a primitive oriented
    determinant-preserving carrier".

\item \textbf{Interface table.}\ A new
    \S\ref{sec:interface-table} maps every load-bearing TFPT
    Paper 2 claim to its current Lean status and the missing
    upstream proof. This separates "formalised", "interfaced",
    and "out of scope" without ambiguity.

\item \textbf{Claim Contract refinement.}\ The contract now
    distinguishes \emph{direct} inputs to
    \texttt{CarrierPremises} (the rank witnesses and the
    boundary involution) from \emph{refinable} upstream
    interfaces (Stages B/C/D feeding the Yukawa rank witness).

\item \textbf{Editorial fixes.}\ The running header is
    overridden for this note (it previously inherited the
    TFPT 4.5 series header). The "v3" relic in the Result
    section is corrected.
\end{itemize}

\subsection*{v6 --- Yukawa Stage D and Calderón interface}

Relative to v5, this version adds two new modules and resolves
the standing review-level critique that "Stage C is data, not
derivation":

\begin{itemize}[leftmargin=1.6em, itemsep=2pt]
\item Added module \texttt{YukawaTrilinearForm.lean} with the
    structure
    \texttt{YukawaTrilinearForm K E := \{$\omega : \Lambda^3 E \to_{\ell K} K$\}}
    and two \texttt{Prop} predicates (originally named
    \texttt{Nondegenerate} and \texttt{Primitive}, renamed in
    v7 to \texttt{ContractionInjective} and
    \texttt{ContractionSurjective} respectively).
\item Constructed the contraction map
    \texttt{contraction} $: E \to_{\ell K} ((\Lambda^2 E) \to_{\ell K} K)$ from $\omega$
    explicitly using
    \texttt{exteriorPower.alternatingMapLinearEquiv} together
    with \texttt{AlternatingMap.curryLeft}. The contraction is
    now a Lean \emph{definition}, not a hypothesis.
\item Proved
    \texttt{contractionEquivOfBoth}: under both predicates the
    contraction upgrades to a linear isomorphism. Hence
    \texttt{toPrimitiveIndecomposableYukawaCoupling}: Stage D
    induces Stage C as a theorem.
\item Proved
    \texttt{finrank\_E\_eq\_three}: Stage D directly implies
    $\dim E = 3$ through the full tower.
\item Added module \texttt{CalderonInterface.lean} with the
    structure \texttt{CalderonCertificate} that wraps Paper 1's
    Calder\'on-induced carrier involution as a Lean datum and
    feeds it into the existing \texttt{CarrierInvolution}
    pipeline. The Paper 1 input is now an explicit
    interface, not a silent consumption.
\item Updated \texttt{AxiomCheck.lean} and
    \texttt{AuditCheck.lean} with eight new theorem names.
    All eight depend only on
    $\{\texttt{propext},\ \texttt{Classical.choice},\
    \texttt{Quot.sound}\}$. \texttt{audit.sh} returns
    \texttt{AUDIT:\ PASS}.
\end{itemize}

\smallskip
The Yukawa interface is now a four-stage tower (under
\texttt{Nontrivial}\,$E_-$ at the Stage D entry point)

\[
\underbrace{(\omega, \texttt{CI}, \texttt{CS})}_{\text{Stage D}}
\;\Rightarrow\;
\underbrace{C_\omega : E_- \simeq (\Lambda^2 E_-)^*}_{\text{Stage C}}
\;\Rightarrow\;
\underbrace{\Lambda^3 E_- \simeq K}_{\text{Stage B}}
\;\Rightarrow\;
\underbrace{E_- \simeq K^3}_{\text{Stage A}}
\;\Rightarrow\;
\dim E_- = 3.
\]

Every step is a Lean theorem. The primary algebraic content
sits at Stage D: an alternating trilinear form plus two
explicit structural predicates.

\subsection*{v5 --- Yukawa Stage C}

Relative to v4, this version closes the Yukawa interface
\emph{within} Lean by introducing a third stage:

\begin{itemize}[leftmargin=1.6em, itemsep=2pt]
\item Added module \texttt{YukawaPrimitive.lean} with the
    structure
    \texttt{PrimitiveIndecomposableYukawaCoupling K E := \{contraction : E $\simeq_{\ell K}$ ($\Lambda^2 E$)\textsuperscript{*}, nontrivial : Nontrivial E\}}.
\item Proved \texttt{finrank\_E\_eq\_three} via the dimension
    chain
    $\dim E \overset{\text{iso}}{=} \dim (\Lambda^2 E)^*
    \overset{\text{dual}}{=} \dim \Lambda^2 E
    \overset{\text{ext.\,pow.}}{=} \binom{\dim E}{2}$
    combined with the elementary fact that $n = n(n-1)/2$ and
    $n \ge 1$ force $n = 3$.
\item Proved the converse \texttt{ofFinrankEqThree} and the
    characterisation \texttt{nonempty\_iff\_finrank\_eq\_three}.
\item Bridged Stage C to Stage B via \texttt{toYukawaTopForm}, and
    transitively to Stage A via \texttt{toYukawaRankCertificate}.
\item Updated \texttt{AxiomCheck.lean} and \texttt{AuditCheck.lean}
    with the five new theorem names. All five depend only on
    $\{\texttt{propext},\ \texttt{Classical.choice},\
    \texttt{Quot.sound}\}$.
\end{itemize}

\smallskip
The Yukawa interface is now structured as a three-stage tower

\[
\underbrace{C_\omega : E_- \simeq (\Lambda^2 E_-)^*}_{\text{Stage C}}
\;\Rightarrow\;
\underbrace{\Lambda^3 E_- \simeq K}_{\text{Stage B}}
\;\Rightarrow\;
\underbrace{E_- \simeq K^3}_{\text{Stage A}}
\;\Rightarrow\;
\dim E_- = 3,
\]

with each step a Lean theorem and a characterisation establishing
the equivalence to $\dim E_- = 3$ at every stage. The
\emph{conceptual} input to Stage C is a contraction iso of $E_-$
with $(\Lambda^2 E_-)^*$, which is the mathematical core of
"primitive indecomposable Yukawa type" with no physics words.

\subsection*{v4 --- Yukawa Stage B}

In v4 the Yukawa rank discharge was promoted from a bare Stage-A
rank witness ($E_- \simeq K^3$, a numerical assumption in
linear-equivalence dress) to a Stage-B algebraic hypothesis
($\Lambda^3 E_- \simeq K$, a structural property of $E_-$).

\begin{itemize}[leftmargin=1.6em, itemsep=2pt]
\item Added module \texttt{YukawaTopForm.lean} with the structure
    \texttt{YukawaTopForm K E := \{topFormEquiv : $\Lambda^3$ E
    $\simeq_{\ell K}$ K\}}.
\item Proved \texttt{YukawaTopForm.finrank\_E\_eq\_three} using
    Mathlib's \texttt{exteriorPower.finrank\_eq} and
    \texttt{Nat.choose\_eq\_one\_iff}.
\item Proved the converse \texttt{ofFinrankEqThree} and the
    characterisation \texttt{nonempty\_iff\_finrank\_eq\_three}.
\item Bridged Stage B to Stage A via
    \texttt{toYukawaRankCertificate}.
\item Updated \texttt{AxiomCheck.lean} and \texttt{AuditCheck.lean}
    with the five new theorem names. All five depend only on
    $\{\texttt{propext},\ \texttt{Classical.choice},\
    \texttt{Quot.sound}\}$.
\end{itemize}

After v4, the conceptual gap shrunk to
\[
\text{"primitive indecomposable Yukawa type"}
\;\Longrightarrow\;
\text{Nonempty}\,(\texttt{YukawaTopForm}\,K\,E_-).
\]
v5 closes this gap with the Stage C contraction iso definition.

\subsection*{v3 --- correcting overclaims}

Relative to v2, v3 corrected three classes of overclaim:

\begin{enumerate}[leftmargin=1.6em, itemsep=3pt]
\item \textbf{Scope of the title and the abstract.} "SM
    Hypercharge Spectrum" $\to$ "$3+2$ Hypercharge Carrier". The
    note no longer suggests it has formalised all SM chiral-
    multiplet hypercharges; only the carrier multiplet
    $\mathrm{diag}(-\tfrac13,-\tfrac13,-\tfrac13,\tfrac12,\tfrac12)$.
\item \textbf{Status of the Stage-A certificates.} "Not
    assumptions" $\to$ "structured Stage-A hypotheses, not bare
    numerical assumptions". The Higgs and Yukawa certificates
    are hypotheses in Lean. Calling them not-assumptions
    misrepresents the type-theoretic status. The Yukawa
    structure is renamed from \texttt{PrimitiveYukawaTypeCertificate}
    to \texttt{YukawaRankCertificate}, accurately describing what
    it actually carries.
\item \textbf{Logical direction of the dependency graph.} The v2
    graph had three incorrect arrows. v3 routes
    \texttt{DeterminantCharacter} $\to$
    \texttt{LatticeRigidityGeneral} (not the reverse); removes
    the spurious \texttt{HiggsIndexCertificate} $\to$
    \texttt{YukawaRankCertificate} edge; and removes the
    \texttt{InvolutionProjectors} $\to$
    \texttt{LatticeRigidityGeneral} edge.
\end{enumerate}

Two further substantive additions:

\begin{itemize}[leftmargin=1.6em, itemsep=1pt]
\item \textbf{Bundled main theorems strengthened.} v3 adds
    \texttt{CarrierPremises.primitive\_pair\_eq\_sm},
    \texttt{CarrierPremises.carrier\_polynomial\_Y}, and the
    top-level \texttt{hypercharge\_carrier\_packet} bundling all
    three conclusions $(q_-,q_+)=(-2,3),\ \tr Y=0,\
    6 Y^2 - Y - \mathbf 1 = \mathbf 0$. The previous version only
    bundled $\tr Y = 0$, which was weaker than the abstract
    claimed.
\item \textbf{Universal-$\lambda$ determinant character.} v3 adds
    \texttt{trace\_zero\_of\_det\_one\_forall\_rat} as the proper
    headline, with \texttt{trace\_zero\_of\_det\_one\_at\_two} kept
    as a single-point helper.
\end{itemize}

The v2-level claim now reduces to its accurate form, refined
for v9:
\[
\boxed{
\begin{aligned}
&\text{Lean verifies the algebraic carrier graph from:}\\
&\quad\text{(i) a Calder\'on \emph{idempotent}
     ($\texttt{CalderonProjector}$: $\pi^2 = \pi$ $\Rightarrow$
     $\varepsilon^2 = 1$),}\\
&\quad\text{(ii) a Higgs Stage-B \emph{polynomial structure}
     ($E_+ \simeq \alpha X + \beta Y$ $\Rightarrow$
     $\dim E_+ = 2$),}\\
&\quad\text{(iii) a Yukawa Stage-D \emph{trilinear form}
     ($\omega : \Lambda^3 E_- \to K$ with
     \texttt{ContractionInjective},}\\
&\quad\text{\phantom{(iii) }\texttt{ContractionSurjective}, and
     \texttt{Nontrivial}\,$E_-$ $\Rightarrow$ $\dim E_- = 3$),}\\
&\quad\text{(iv) a primitive oriented
     determinant-preserving carrier condition,}\\
&\text{down to the $3+2$ hypercharge carrier
     $Y = \mathrm{diag}(-\tfrac13,-\tfrac13,-\tfrac13,\tfrac12,\tfrac12)$}\\
&\text{with $(q_-, q_+) = (-2, 3)$, $\tr Y = 0$,
     and $6 Y^2 - Y - \mathbf 1 = \mathbf 0$.}
\end{aligned}
}
\]
All three rank/involution discharges are now \emph{structural}:
the Calder\'on input is an algebraic idempotent (not a wrapped
involution), the Higgs input is a polynomial-space iso (not a
bare rank witness), and the Yukawa input is an alternating
trilinear form plus two operational predicates (not a packaged
iso).

% ---------------------------------------------------------------- %
\section{Next Steps}
\label{sec:next}
% ---------------------------------------------------------------- %

In rough order of ascending intellectual difficulty:

\begin{enumerate}[leftmargin=1.6em, itemsep=2pt]
\item \textbf{CI workflow.}
    Wire \texttt{audit.sh} into a GitHub Actions workflow that
    refuses any merge with \texttt{AUDIT:\ FAIL}.

\item \textbf{Higgs certificate from $\mathbb P^1$ cohomology.}
    Build a Lean construction of \texttt{HiggsIndexCertificate}
    from the cohomology of $\mathcal O(1)$ on $\mathbb P^1$.

\item \textbf{Higgs scheme-cohomology realisation.}\ The
    Stage B \texttt{HiggsTopForm} now identifies $E_+$ with
    the algebraic shadow of $H^0(\mathbb P^1, \mathcal O(1))$;
    a future module should build the scheme-theoretic
    $H^0(\mathbb P^1_K, \mathcal O(1))$ and exhibit an iso to
    \texttt{DegreeOneTwoVarForm K}. This is Mathlib algebraic
    geometry work and reasonably sized.
\item \textbf{Calderón boundary spectral pair.}\ The Stage A+
    \texttt{CalderonProjector} derives $\varepsilon^2 = 1$
    from an algebraic idempotent, and v10's
    \texttt{BoundaryPolarization} derives the idempotent from
    a complement-pair $H = H_+ \oplus H_-$. What remains is the
    analytic Paper-1 theorem producing the polarisation itself
    from boundary elliptic theory (Calder\'on operator on
    Sobolev spaces). This is a major independent
    formalisation project.
\item \textbf{Yukawa: existence of $\omega$ from TFPT
    \emph{canonically}.}\ The Lean chain now provides full
    existence:
    $\dim E_- = 3 \Rightarrow \exists\,\omega$ with
    \texttt{ContractionInjective} and
    \texttt{ContractionSurjective}
    (\texttt{YukawaStageDExistence}, v10). What remains is
    the TFPT-internal question: does the boundary kernel
    produce \emph{the same canonical} $\omega$ (or one
    equivalent up to the action of the carrier symmetries)?
    The abstract bridge
    \texttt{BoundaryYukawaKernelInterface} (v11) gives a
    Lean-typed target for this future TFPT theorem.

\item \textbf{Derived SM representations.}
    Implement
    \texttt{LinearMap.dualHypercharge},
    \texttt{ExteriorPower.hypercharge},
    \texttt{TensorProduct.hypercharge},
    so that the project can claim "SM hypercharge spectrum"
    rather than "carrier hypercharge".

\item \textbf{$\mathbb Z_6$-quotient theorem.}
    The kernel of the carrier action of
    $SU(3)\times SU(2)\times U(1)$ on $\mathbb K^5$ as $\mathbb Z_6$,
    using \texttt{Mathlib.LinearAlgebra.UnitaryGroup}.

\item \textbf{Mathlib upstream.}
    After (2) and (4), \texttt{Polarization},
    \texttt{InvolutionProjectors.toPolarization}, and the
    Mathlib bridge make a clean PR candidate.
\end{enumerate}

% ---------------------------------------------------------------- %
\section{Result}
% ---------------------------------------------------------------- %

\begin{tfptexports}
\textbf{Lean-verified outputs (v11).}
\begin{itemize}[leftmargin=1.5em, itemsep=0pt]
\item Abstract carrier polynomial:
  \texttt{Polarization.sixY\_carrier\_polynomial}.
\item Boundary involution $\Rightarrow$ projectors:
  \texttt{InvolutionProjectors.CarrierInvolution.toPolarization}.
\item Mathlib bridge:
  \texttt{MathlibBridge.Polarization.toCompleteOrthogonalIdempotents}.
\item General integer rigidity:
  \texttt{LatticeRigidityGeneral.primitive\_trace\_free\_pair\_general}.
\item Structural trace identity:
  \texttt{TraceProjection.trace\_linear\_combination\_of\_idempotents}.
\item Determinant character, universal-$\lambda$ headline:
  \texttt{DeterminantCharacter.trace\_zero\_of\_det\_one\_forall\_rat}.
\item Higgs Stage-A discharge:
  \texttt{HiggsIndexShadow.HiggsIndexCertificate.finrank\_Eplus\_eq\_two}.
\item Yukawa Stage-A discharge (rank witness):
  \texttt{YukawaRank.YukawaRankCertificate.finrank\_Eminus\_eq\_three}.
\item Yukawa Stage-B discharge (algebraic top form):
  \texttt{YukawaTopForm.YukawaTopForm.finrank\_E\_eq\_three}.
\item Yukawa Stage-B characterisation:
  \texttt{YukawaTopForm.YukawaTopForm.nonempty\_iff\_finrank\_eq\_three}.
\item Stage-B $\to$ Stage-A bridge:
  \texttt{YukawaTopForm.YukawaTopForm.toYukawaRankCertificate}.
\item Yukawa Stage-C discharge (contraction iso):
  \texttt{YukawaPrimitive.PrimitiveIndecomposableYukawaCoupling.finrank\_E\_eq\_three}.
\item Yukawa Stage-C characterisation:
  \texttt{YukawaPrimitive.PrimitiveIndecomposableYukawaCoupling.nonempty\_iff\_finrank\_eq\_three}.
\item Stage-C $\to$ Stage-B bridge:
  \texttt{YukawaPrimitive.PrimitiveIndecomposableYukawaCoupling.toYukawaTopForm}.
\item Yukawa Stage-D primary datum and constructed contraction:
  \texttt{YukawaTrilinearForm.YukawaTrilinearForm.contraction}.
\item Stage-D operational predicates:
  \texttt{YukawaTrilinearForm.YukawaTrilinearForm.ContractionInjective} and \texttt{.ContractionSurjective}.
\item Stage-D combined predicate
  (\texttt{PrimitiveIndecomposableYukawaCondition}) and
  rank-3 discharge under \texttt{Nontrivial}\,$E$:
  \texttt{YukawaTrilinearForm.YukawaTrilinearForm.finrank\_E\_eq\_three}
  and \texttt{.finrank\_E\_eq\_three\_of\_condition}.
\item Stage-D transparency form (no \texttt{Nontrivial}, $\dim
  \in \{0,3\}$): \texttt{YukawaTrilinearForm.YukawaTrilinearForm.finrank\_E\_eq\_zero\_or\_three}.
\item Stage-D $\to$ Stage-C bridge:
  \texttt{YukawaTrilinearForm.YukawaTrilinearForm.toPrimitiveIndecomposableYukawaCoupling}.
\item Calderón certificate interface:
  \texttt{CalderonInterface.CalderonCertificate.toCarrierInvolution}.
\item Calderón structural derivation (idempotent $\Rightarrow$
  involution): \texttt{CalderonProjector.CalderonProjector.eps\_sq}.
\item Calderón upstream (v10): boundary polarisation $\Rightarrow$
  idempotent:
  \texttt{BoundaryPolarization.BoundaryPolarization.projector\_isIdempotent}
  and \texttt{toCalderonProjector}.
\item Higgs Stage B discharge (polynomial space $\Rightarrow$
  $\dim E_+ = 2$):
  \texttt{HiggsTopForm.HiggsTopForm.finrank\_Eplus\_eq\_two}.
\item Higgs scheme-cohomology shadow (v10):
  \texttt{HiggsSchemeCohomologyShadow.toDegreeOneTwoVarForm}
  (\texttt{LinearEquiv} between the algebraic shadow and
  \texttt{MvPolynomial.homogeneousSubmodule (Fin 2) K 1}).
\item Yukawa Stage D backward (v10):
  \texttt{YukawaStageDExistence.exists\_yukawaTrilinearForm\_of\_finrank\_eq\_three}
  and the characterisation
  \texttt{stage\_D\_iff\_finrank\_eq\_three} (under
  \texttt{Nontrivial}\,$E$).
\item Boundary polarisation converse (v11):
  \texttt{BoundaryPolarization.ofCarrierInvolution}
  --- the Paper-1 interface is a Lean-proved \emph{algebraic}
  equivalence between the polarisation picture and the
  involution picture (two presentations of the same
  finite-dimensional carrier datum; not the analytic
  construction).
\item Seam-winding interface (v11):
  \texttt{SeamWindingInterface.SeamWindingData.toPrimitiveOrientedDeterminantCarrier}
  and the converse, with the SM specialisation
  \texttt{sm\_pair\_eq}.
\item Boundary Yukawa kernel interface (v11):
  \texttt{BoundaryYukawaKernelInterface.BoundaryYukawaKernel.finrank\_E\_eq\_three},
  \texttt{ofFinrankEqThree}, and the characterisation
  \texttt{nonempty\_iff\_finrank\_eq\_three}. \emph{Typed
  packaging of the Stage D condition --- not a new physical
  derivation. The canonicity of $\omega$ from the TFPT
  boundary kernel remains upstream physics work.}
\item Bundled main theorems:
  \texttt{CarrierPremises.primitive\_pair\_eq\_sm},
  \texttt{CarrierPremises.trace\_Y\_eq\_zero},
  \texttt{CarrierPremises.carrier\_polynomial\_Y},
  \texttt{CarrierPremises.hypercharge\_carrier\_packet}.
\item Concrete carrier polynomial:
  \texttt{Hypercharge.Y\_carrier\_polynomial}.
\end{itemize}

\textbf{Audit outputs.}
$1788$ build jobs green; zero \texttt{sorry}; zero
\texttt{admit}; zero \texttt{unsafe} / \texttt{opaque} /
\texttt{partial}; zero domain-specific axioms; every headline
theorem depends only on
$\{\texttt{propext},\ \texttt{Classical.choice},\ \texttt{Quot.sound}\}$.
\end{tfptexports}

% ---------------------------------------------------------------- %
\appendix
\section{Reproducing the Build and Audit}
% ---------------------------------------------------------------- %

\begin{lstlisting}[language=bash]
# 1. Install elan (Lean toolchain manager)
curl -sSf https://raw.githubusercontent.com/leanprover/elan/master/elan-init.sh \
  | sh -s -- -y --default-toolchain none
export PATH="$HOME/.elan/bin:$PATH"

# 2. Enter the project
cd experiments/lean4-carrier-rigidity

# 3. Fetch Mathlib and its pre-built oleans
lake update                  # ~30 s
lake exe cache get           # ~3 min, ~1 GB

# 4. Build the project
lake build                   # ~8 s on top of cache; 1788 jobs

# 5. Run the seven-check audit gate
./scripts/audit.sh
\end{lstlisting}

Expected trailer:

\begin{lstlisting}
=== Final verdict ===
  AUDIT: PASS
\end{lstlisting}

\end{document}
